%%
%% This is file `sample-sigconf-authordraft.tex',
%% generated with the docstrip utility.
%%
%% The original source files were:
%%
%% samples.dtx  (with options: `all,proceedings,bibtex,authordraft')
%% 
%% IMPORTANT NOTICE:
%% 
%% For the copyright see the source file.
%% 
%% Any modified versions of this file must be renamed
%% with new filenames distinct from sample-sigconf-authordraft.tex.
%% 
%% For distribution of the original source see the terms
%% for copying and modification in the file samples.dtx.
%% 
%% This generated file may be distributed as long as the
%% original source files, as listed above, are part of the
%% same distribution. (The sources need not necessarily be
%% in the same archive or directory.)
%%
%%
%% Commands for TeXCount
%TC:macro ~\cite [option:text,text]
%TC:macro ~\citep [option:text,text]
%TC:macro ~\citet [option:text,text]
%TC:envir table 0 1
%TC:envir table* 0 1
%TC:envir tabular [ignore] word
%TC:envir displaymath 0 word
%TC:envir math 0 word
%TC:envir comment 0 0
%%
%% The first command in your LaTeX source must be the \documentclass
%% command.
%%
%% For submission and review of your manuscript please change the
%% command to \documentclass[manuscript, screen, review]{acmart}.
%%
%% When submitting camera ready or to TAPS, please change the command
%% to \documentclass[sigconf]{acmart} or whichever template is required
%% for your publication.
%%
%%
% \documentclass[sigconf]{acmart}
% \documentclass[manuscript, screen, review]{acmart}
\documentclass[sigconf]{acmart}

% \usepackage{cite}
% \usepackage{amsmath,amssymb,amsfonts}

\usepackage{graphicx}
\usepackage{textcomp}
\usepackage{xcolor}
% \usepackage[a4paper, total={184mm,239mm}]{geometry}
\usepackage{subcaption}
\usepackage{svg}
\usepackage{float}
\usepackage{multicol}
\usepackage{array}
\usepackage{mwe}
\usepackage[noend]{algpseudocode}
\usepackage[ruled, lined, longend]{algorithm2e}
\usepackage{multirow}
\usepackage{arydshln}
\usepackage{adjustbox}
\usepackage[thinc]{esdiff}
\usepackage{booktabs} % To thicken table lines
\usepackage{pifont}% http://ctan.org/pkg/pifont
% \usepackage{mathabx}
\usepackage{hyperref}
\usepackage{todonotes}
\usepackage{enumitem}
\usepackage{tablefootnote}
\usepackage{threeparttable}
\usepackage{booktabs}
\usepackage{multirow}
\usepackage{amsmath}
\usepackage{siunitx}
\usepackage{graphicx}
\newtheorem{remark}{Remark}

%%
%% \BibTeX command to typeset BibTeX logo in the docs
\AtBeginDocument{%
  \providecommand\BibTeX{{%
    Bib\TeX}}}

%% Rights management information.  This information is sent to you
%% when you complete the rights form.  These commands have SAMPLE
%% values in them; it is your responsibility as an author to replace
%% the commands and values with those provided to you when you
%% complete the rights form.
\setcopyright{none}
% \setcopyright{acmlicensed}
% \copyrightyear{2018}
% \acmYear{2018}
% \acmDOI{XXXXXXX.XXXXXXX}
%% These commands are for a PROCEEDINGS abstract or paper.
\acmConference[]{}
%%
%%  Uncomment \acmBooktitle if the title of the proceedings is different
%%  from ``Proceedings of ...''!
%%
%%\acmBooktitle{Woodstock '18: ACM Symposium on Neural Gaze Detection,
%%  June 03--05, 2018, Woodstock, NY}
%% \acmISBN{978-1-4503-XXXX-X/2018/06}


%%
%% Submission ID.
%% Use this when submitting an article to a sponsored event. You'll
%% receive a unique submission ID from the organizers
%% of the event, and this ID should be used as the parameter to this command.
%%\acmSubmissionID{123-A56-BU3}

%%
%% For managing citations, it is recommended to use bibliography
%% files in BibTeX format.
%%
%% You can then either use BibTeX with the ACM-Reference-Format style,
%% or BibLaTeX with the acmnumeric or acmauthoryear sytles, that include
%% support for advanced citation of software artefact from the
%% biblatex-software package, also separately available on CTAN.
%%
%% Look at the sample-*-biblatex.tex files for templates showcasing
%% the biblatex styles.
%%

%%
%% The majority of ACM publications use numbered citations and
%% references.  The command ~\citestyle{authoryear} switches to the
%% "author year" style.
%%
%% If you are preparing content for an event
%% sponsored by ACM SIGGRAPH, you must use the "author year" style of
%% citations and references.
%% Uncommenting
%% the next command will enable that style.
%%~\citestyle{acmauthoryear}


%%
%% end of the preamble, start of the body of the document source.

\settopmatter{printacmref=false}

\begin{document}

%%
%% The "title" command has an optional parameter,
%% allowing the author to define a "short title" to be used in page headers.

% \title{TeLLMe: An Energy-Efficient \underline{Te}rnary \underline{LLM} Accelerator for Prefill and Decode on \underline{E}dge FPGAs}

%% \title{TeLLMe: A \underline{Te}rnary Table-Lookup Based Accelerator Enabling End-to-End \underline{LLM} Prefill and Decode on \underline{E}dge FPGAs}

\title{TeLLMe v2: An Efficient End-to-End \underline{Te}rnary \underline{LLM} Prefill and Decode Accelerator with Table-Lookup Matmul on \underline{E}dge FPGAs}
\author{Ye Qiao}
\authornote{Equal contribution.}
\email{yeq6@uci.edu}
\affiliation{%
  \institution{University of California, Irvine}
  \city{Irvine}
  \state{California}
  \country{USA}
}

\author{Zhiheng Chen}
\authornotemark[1] % reuse the first author’s note mark
\email{zhihenc5@uci.edu}
\affiliation{%
  \institution{University of California, Irvine}
  \city{Irvine}
  \state{California}
  \country{USA}
}

\author{Yifan Zhang}
\email{yifanz58@uci.edu}
\affiliation{%
  \institution{University of California, Irvine}
  \city{Irvine}
  \state{California}
  \country{USA}
}

\author{Yian Wang}
\email{yianw11@uci.edu}
\affiliation{%
  \institution{University of California, Irvine}
  \city{Irvine}
  \state{California}
  \country{USA}
}

\author{Sitao Huang}
\email{sitaoh@uci.edu}
\affiliation{%
  \institution{University of California, Irvine}
  \city{Irvine}
  \state{California}
  \country{USA}
}

% \author{Zhiheng Chen\textsuperscript{\S}} 
% \author{Yifan Zhang} 
% \author{Yian Wang} 
% \author{Sitao Huang}
%     Ye Qiao\textsuperscript{\S}, Zhiheng Chen\textsuperscript{\S}, Yifan Zhang, Yian Wang, Sitao Huang \\
%     \textit{Department of Electrical Engineering and Computer Science} \\
%     \textit{University of California,} Irvine, USA \\
% \email{yeq6, zhihenc5, yifanz58, yianw11, sitaoh\}@uci.edu
% \authornote{\textsuperscript{\S}Equal contribution}
%%
%% The "author" command and its associated commands are used to define
%% the authors and their affiliations.
%% Of note is the shared affiliation of the first two authors, and the
%% "authornote" and "authornotemark" commands
%% used to denote shared contribution to the research.
% \author{Ben Trovato}
% \authornote{Both authors contributed equally to this research.}
% \email{trovato@corporation.com}
% \orcid{1234-5678-9012}
% \author{G.k.m. Tobin}
% \authornotemark[1]
% \email{webmaster@marysville-ohio.com}
% \affiliation{%
%   \institution{Institute for Clarity in Documentation}
%   \city{Dublin}
%   \state{Ohio}
%   \country{USA}
% }

% \author{Lars Th{\o}rv{\"a}ld}
% \affiliation{%
%   \institution{The Th{\o}rv{\"a}ld Group}
%   \city{Hekla}
%   \country{Iceland}}
% \email{larst@affiliation.org}

% \author{Valerie B_{\text{idx}}\'eranger}
% \affiliation{%
%   \institution{Inria Paris-Rocquencourt}
%   \city{Rocquencourt}
%   \country{France}
% }

% \author{Aparna Patel}
% \affiliation{%
%  \institution{Rajiv Gandhi University}
%  \city{Doimukh}
%  \state{Arunachal Pradesh}
%  \country{India}}

% \author{Huifen Chan}
% \affiliation{%
%   \institution{Tsinghua University}
%   \city{Haidian Qu}
%   \state{Beijing Shi}
%   \country{China}}

% \author{Charles Palmer}
% \affiliation{%
%   \institution{Palmer Research Laboratories}
%   \city{San Antonio}
%   \state{Texas}
%   \country{USA}}
% \email{cpalmer@prl.com}

% \author{John Smith}
% \affiliation{%
%   \institution{The Th{\o}rv{\"a}ld Group}
%   \city{Hekla}
%   \country{Iceland}}
% \email{jsmith@affiliation.org}

% \author{Julius P. Kumquat}
% \affiliation{%
%   \institution{The Kumquat Consortium}
%   \city{New York}
%   \country{USA}}
% \email{jpkumquat@consortium.net}

%%
%% By default, the full list of authors will be used in the page
%% headers. Often, this list is too long, and will overlap
%% other information printed in the page headers. This command allows
%% the author to define a more concise list
%% of authors' names for this purpose.

% \renewcommand{\shortauthors}{Trovato et al.}

%%
%% The abstract is a short summary of the work to be presented in the
%% article.

\begin{abstract}
%% needs to explain why it is necessary to deploy LLMs on edge devices.
%\vspace{-1mm}
With the emergence of wearable devices and other embedded systems, 
deploying large language models (LLMs) on edge platforms becomes an urgent need. However, it is challenging because of their high computational and memory demands. Although recent low-bit quantization methods (e.g., BitNet, DeepSeek) compress weights to as low as 1.58 bits with minimal accuracy loss, edge deployment is still constrained by limited on-chip resources, power budgets, and the often-neglected long latency of the prefill stage. 
%% However, while table-lookup-based methods have been explored for perceptron and convolutional networks on FPGAs, their application to LLMs on FPGAs remains unexplored. 
We present TeLLMe, the first table-lookup-based ternary LLM accelerator for low-power edge FPGAs that fully supports both prefill and autoregressive decoding using 1.58-bit weights and 8-bit activations. TeLLMe incorporates our proposed novel techniques including (1) a table-lookup-based ternary matrix multiplication (TLMM) engine utilizing grouped activations and online precomputation for low resource utilization and high throughput; (2) a fine-grained analysis and optimization onan analytic and fine-grained URAM-based weight buffer management of weight loading from global memory and compute engine weight access; (3) a streaming dataflow architecture that fuses floating-point element-wise operations with linear computations to hide latency; (4) a reversed-reordered prefill stage attention with fused attention operation for high memory efficiency; and (5) a resource-efficient specialized decoding stage attention. Under a 5W power budget, TeLLMe delivers up to 25 tokens/s decoding throughput and 0.45s to 0.96s Time-to-First-Token (TTFT) for 64–128 token prompts, marking a significant energy-efficiency advancement in LLM inference on edge FPGAs. %%  and establishing a new edge FPGA benchmark for generative AI.
% Deploying large language models (LLMs) on edge
%  platforms is challenged by their high computational and memory
%  demands. Although recent low-bit quantization methods (e.g.,
%  BitNet, DeepSeek) compress weights to as little as 1.58 bits with
%  minimal accuracy loss, edge deployment is still constrained by
%  limited on-chip resources, power budgets, and the often-neglected
%  long latency of the prefill phase. We present TeLLMe, the first ternary
%  LLM accelerator for low-power FPGAs (e.g., AMD KV260)
%  that fully supports both prefill and autoregressive decoding
%  using 1.58-bit weights and 8-bit activations. Our contributions
%  include: (1) a table-lookup-based matrix multiplication (TLMM) inspired ternary matrix multiplication engine utilizing grouped activations and online precomputation for low resource utilization and high throughput looking up; (2) an analytic and fine-grained URAM weight buffer management of weight loading from DDR and compute engine weight access; (3) a streaming dataflow architecture that fuses floating-point element-wise operations with linear computations to hide latency ;(4)a reversed-reordered prefill stage attention with fused QKV operation and a on-chip memory-supported specialized decoding stage attention. 
%  Under a 5W power budget, TeLLMe delivers up to 25 token/s
%  decoding throughput and prefill latencies of 0.46s to 0.98s for 64–128 token prompts, marking a significant
%  energy-efficiency advance and establishing a new edge FPGA
%  benchmark for generative AI.
 
% Deploying large language models (LLMs) on edge platforms poses significant challenges due to their substantial computational and memory requirements. Recent advances in low-bit quantization, such as BitNet and DeepSeek, have demonstrated that LLM weights can be compressed to as low as 1.58 bits without significant accuracy degradation. However, practical edge deployment remains limited by tight on-chip resource constraints and power budgets. Notably, prior works often overlook the substantial latency introduced during the prefill stage of LLM inference, which is critical for real-time and safety-sensitive applications. In this paper, we introduce \textbf{TeLLMe}, the first ternary LLM accelerator for edge FPGAs with full support for both the prefill and autoregressive decoding phases. Targeting low-power platforms such as the AMD KV260, TeLLMe supports 1.58-bit quantized weights and 8-bit activations. Our design features: (1) Table-lookup-based matrix multiplication (TLMM) inspired ternary matrix multiplication engine utilizing grouped activations and online precomputation for low resource utilization and high throughput; (2) Analytic parameter selection of TLMM and fine-grained URAM weight buffer management of DDR weight loading and compute engine weight access; (3) Streaming dataflow fusion of floating point element-wise operation and TLMM for latency overlapping ;(4)a reversed-reordered prefill stage attention with fused QKV operation and a on-chip memory-supported specialized decoding stage attention. 
% TeLLMe achieves up to 25 tokens/s decoding throughput and supports 1024-token context lengths, all while operating under 5 watts. TeLLMe also achieves prefill latencies ranging from 0.46s to 0.98s for prompt sizes of 64 to 128 tokens, demonstrating practical viability and deployment potential in fast-response real-world applications. These results demonstrate a significant improvement in energy efficiency over existing edge FPGA and edge GPU accelerators. To our knowledge, TeLLMe is the first edge FPGA accelerator to combine ternary quantization with comprehensive prefill and decoding support, setting a new benchmark for energy-efficient generative AI at the edge.

% This represents a substantial energy-efficiency improvement over current edge FPGA accelerators. To the best of our knowledge, TeLLMe is the first edge FPGA-based LLM accelerator to combine ternary quantization with both prefill and decoding inference support, setting a new benchmark for energy-efficient generative AI deployment at the edge.
\end{abstract}



\maketitle

\section{Introduction}
%\vspace{-1mm}
Large language models (LLMs) have been evolving rapidly, delivering state-of-the-art results in machine translation, code generation, question answering, and conversational AI. Models such as GPT-3~\cite{brown2020language}, LLaMA~\cite{touvron2023LLaMA}, and DeepSeek-R1~\cite{guo2025deepseek} highlight the benefits of scale, but those gains arrive with steep cost increases in computation, memory, and energy.

Deploying LLMs on edge devices (embedded CPUs, GPUs, FPGAs, etc.) preserves privacy, reduces latency, and enables autonomy, yet it is difficult due to tight limits on memory bandwidth and capacity, compute resources, and power budgets. Autoregressive decoding further stresses these constraints: growing key–value (KV) caches, long-context handling, and strict latency requirements often become the dominating performance bottlenecks. %%  the on-device critical path.

Extreme model compression via low-bit quantization has emerged as a promising technique~\cite{qiao2022two,10025006}. BitNet~\cite{bitnet,qiao2025cobra, qiao2025tellme} showed that Transformers can be trained with 1-bit weights; BitNet-1.58~\cite{bitnet158} and DeepSeek~\cite{deepseek} extend this idea to ternary quantization (${-1,0,+1}$), approaching full-precision quality while drastically reducing model size and energy. However, closing the gap between algorithmic compression and efficient, end-to-end deployment on actual edge hardware requires a co-design solution that simultaneously optimizes compute, memory, and scheduling.

While researchers have proposed several solutions to accelerate the decoding stage of LLMs, the existing works often ignore the significant  latency of the prefill stage, making prefill a critical performance bottleneck in the end-to-end LLM systems. 
%% A critical yet often overlooked challenge is the significant prefill latency, which causes the 
%% imbalance between optimizing decoding throughput and optimizing prefill latency. 
For example,~\cite{li2025pushing} demonstrates efficient decoding on embedded FPGAs but leaves the prefill stage largely unaddressed. On device, prefill latency directly impacts users' perceived responsiveness and safety; it is not hidden behind cloud-scale parallelism. 
%% Prefill therefore deserves first-class treatment alongside decoding when designing an edge accelerator.
Therefore, prefill should also be treated as a first-class citizen and carefully accelerated  alongside decoding when designing an edge accelerator. 

%% \footnote{can be pronounced as ``Tell me''}

We present \textbf{TeLLMe}—the \textbf{Te}rnary \textbf{L}arge \textbf{L}anguage \textbf{M}odel \textbf{e}dge accelerator—an edge FPGA accelerator, to our knowledge, the first to incorporate a table-lookup (TL) matrix multiplication (MatMul) approach tailored for ternary LLM inference with full support for both \emph{prefill} stage and \emph{decoding} stage. TeLLMe targets cost-effective low-power FPGAs (AMD Kria KV260), supports 1.58-bit (ternary) weights and 8-bit activations, and co-optimizes compute, memory, and scheduling for low-latency, energy-efficient LLM inference.

Prior FPGA accelerators~\cite{LUTNET, SUMLUTNET, TLMAC} leveraged table-lookup techniques to speed up low-precision arithmetics in other domains, demonstrating the viability of lookup table (LUT)-centric accelerators. However, to our knowledge, TeLLMe is the first to apply a \emph{ternary, table-lookup matmul} design to end-to-end LLM inference (prefill and decoding). Furthermore, it integrates LLM-specific optimizations—fused attention for prefill, disaggregated prefill/decoding dataflows, streaming fusion of dequantization, quantization, element-wise operations, etc., and URAM-aware weight orchestration. With these optimizations, our proposed TeLLMe delivers state-of-the-art efficiency on actual FPGA boards.
Our key contributions include:
\begin{itemize}[leftmargin=*, nosep]
\item \textbf{End-to-end accelerator for ternary LLMs.} We build, to our knowledge, the first edge FPGA accelerator that employs a table-lookup matmul engine for ternary LLMs with full support for both prefill and decoding, achieving up to 25 tokens/s generation throughput and up to 143 tokens/s prefill throughput while consuming under 5~W.
\item \textbf{Ternary table-lookup matmul.} We propose a table-lookup-based matrix multiplication (TLMM) unit optimized for FPGAs that reuses grouped activations and performs online accumulation for ternary multiplications across attention projections and the feed-forward network (FFN).
\item \textbf{URAM-aware weight orchestration.} We introduce fine-grained URAM buffering for high-throughput weight streaming from off-chip memory and provide an analytic method to optimize TLMM engine parameter selection under URAM and LUT constraints.
\item \textbf{Streaming fusion with mixed precision.} We fuse floating-point (FP) dequantization, integer (INT) quantization, and elementwise operations (residual add, rotary embedding, etc)  around the INT-based TLMM to enable mixed-precision execution and overlap compute with dataflow architecture.
\item \textbf{Disaggregated prefill and decoding attention.} We separate the attention pipelines to match their distinct compute/memory patterns. For prefill, we design a fused attention unit with a reversed-attention mechanism and a fully fused pipeline that minimizes off-chip traffic and avoids redundant masked computation. For decoding, we exploit on-chip memory to retain intermediate softmax scores and prevent DDR reloads.
% \item \textbf{Measured efficiency.} On low-power edge hardware, TeLLMe achieves up to 25 tokens/s generation throughput and up to 143 tokens/s prefill throughput while consuming under 5~W.
\end{itemize}
In summary, TeLLMe demonstrates that a LUT/URAM-native, table-lookup matmul engine co-designed with LLM-specific dataflows can unlock end-to-end ternary inference on resource-constrained FPGAs. To our knowledge, it is the first accelerator to realize table-lookup matmul for LLMs with comprehensive support for prefill and decoding on real hardware, establishing a strong baseline for energy-efficient, low-latency generative AI at the edge.

% Large Language Models (LLMs) have achieved remarkable progress in recent years, powering state-of-the-art performance in natural language processing tasks such as machine translation, code generation, question answering, and conversational AI. Models like GPT-3~\cite{brown2020language}, LLaMA~\cite{touvron2023LLaMA}, and DeepSeek-R1~\cite{guo2025deepseek} have shown that increasing model size significantly improves generalization and task performance. However, this scaling comes with substantial costs in terms of computational demands, memory usage, and energy consumption.

% Edge deployment of LLMs, i.e., running these models on low-power, resource-constrained devices such as embedded systems, FPGAs, or mobile SoCs, is a critical enabler for privacy-preserving, latency-sensitive, and autonomous applications. However, such a deployment remains challenging due to the gap between LLM complexity and the limited memory bandwidth, memory capacity, compute capacity, and power budgets on edge platforms.

% To bridge this gap, recent research has focused on \textit{extreme model compression}, particularly through low-bit quantization~\cite{qiao2022two,10025006}. Pioneering work such as BitNet~~\cite{bitnet,qiao2025cobra} demonstrated that Transformer models can be trained with 1-bit weights, while BitNet-1.58~~\cite{bitnet158} and DeepSeek~~\cite{deepseek} extend this to ternary quantization ($\{-1, 0, +1\}$), achieving near-parity model performance with full-precision models. These innovations significantly reduce model size and energy cost, making LLMs more viable for edge execution.

% However, deploying these compressed models on real edge hardware, presents unique challenges. Unlike cloud-scale GPUs, edge FPGAs have strict constraints on on-chip memory (BRAM/URAM), external DDRmatmul approach that specifically designed for ternary LLM inference with full support for both \textit{prefill} and \textit{decoding} stages. TeLLMe enables low-latency, energy-efficient deployment of LLMs on resource-constrained platforms by targeting cost-effective FPGAs such as AMD Kria KV260. It supports ternary-quantized weights (1.58-bit) and 8-bit activations.
% % , delivering scalable inference for models ranging from 700M to 3B parameters.

% Our design co-optimizes compute, memory, and scheduling efficiency. The key contributions of this work can be summarized as follows:
% \begin{itemize}
%     \item We develop the first end-to-end edge FPGA accelerator with table-lookup matmul engine for ternary LLMs supporting both prefill and decoding stages.
%     \item We propose ternary table-lookup-based matrix multiplication (TLMM), an efficient and resource-saving matrix multiplication unit that is specially optimized for FPGA, reusing grouped activations and online computation for ternary matrix multiplications across the projection and feed forward network (FFN).
%     \item We introduce the fine-grained URAM buffer management in high throughput weight loading from DDR and TLMM access and an analytic method to set the TLMM engine parameter within the on-chip URAM and look-up-table unit (LUT) resource limit.
%     \item We introduce the streaming fusion between floating point (FP) dequantization, integer (INT) quantization, and elementwise operation (residual add, mult, and rotary embedding), with the INT-based TLMM module to support mixed-precision and overlap execution.
%     \item We notice the difference of the compute pattern in the prefill and decoding stage attention and disaggregated them into two units. As for the prefill stage, we introduce a fused attention unit, incorporating a novel reversed attention mechanism and a fully fused pipeline to minimize off-chip data movement, avoid redundant masked computation, and guarantee parallelism at the same time. As for the decoding phase, we fully utilize the on-chip memory to buffer the intermediate softmax score matrix to avoid reload it from the dynamic random access memory (DDRmatmul to provide end-to-end support for ternary LLM inference—including both prefill and decoding stages—on real FPGA hardware, establishing a new baseline for energy-efficient, low-latency generative AI at the edge.


% TeLLMe achieves a prefill latency from 0.55s to 1.15s for prompt sizes of 64 to 128 tokens and delivers up to \textbf{9.51 tokens/s} decoding throughput with support for \textbf{1024-token context lengths} on edge FPGAs, all while operating under \textbf{5 watts} of power consumption. This marks a significant advancement over existing mobile-edge devices and prior FPGA-based accelerators, which typically require higher precision and greater power budgets.



%\vspace{-1mm}
\section{Background and Related Work}
%\vspace{-1mm}
% \subsection{Ternary LLM Basic}
% % \listoftodos
% % A typical LLM, such as LLaMA, is composed of multiple identical transformer blocks, each containing an attention module followed by a multilayer perceptron (MLP) module, as illustrated in Fig. \ref{Fig:Bitnet}. Within each attention module, three linear projections are used to compute the query (Q), key (k), and value (V) representations. These are then processed by a multi-head attention mechanism that incorporates both the current QKV tensors and historical KV cache. The MLP module consists of an up-projection and down-projection layer, along with an additional gating projection applied to the up-projection output.

% % The generative inference of LLMs is typically divided into two stages: the \textbf{prefill phase} and the \textbf{decoding phase} (generation), as shown in Fig. \ref{Fig:Bitnet}. During the prefill phase, the entire prompt is processed through the full model stack to produce the first output token. This phase involves parallel computation across multiple input tokens and is dominated by compute-intensive matrix–matrix multiplications, particularly within the linear projection layers. In contrast, the decoding phase proceeds in an autoregressive fashion, generating one token at a time by feeding the previously generated token back into the model. This phase is typically memory-bound due to its reliance on KV cache lookup and smaller matrix–vector operations.

% % Following the observations in Chen et al.~\cite{chen2024understanding} , while FPGAs are generally less efficient than GPUs during the compute-heavy prefill stage, they exhibit competitive advantages during the memory-intensive decoding phase. In this work, we prioritize optimizing both phases of LLM inference to fully leverage the strengths of FPGA architectures.
% The typical BitNet architecture ~\cite{bitnet158} is introduced below and depicted in Fig. \ref{Fig:Bitnet}.
% For a weight matrix $\bf W$, elements are quantized to ${-1, 0, 1}$ via:
% \begin{align}
% {\bf W}_{\text{ternary}} = \text{round}\left(\frac{\bf W}{\alpha}\right),
% \end{align}
% where $\alpha$ is a scaling factor. Matrix multiplication ${\bf Y} = {\bf X}{\bf W}_{\text{ternary}}$ then primarily involves additions and subtractions, since $W_{kj} \in {-1, 0, 1}$:
% \begin{align}
% Y_{ij} = \sum_{k=1}^{d} X_{ik} \cdot W_{kj},
% \end{align}
% For an input activation vector $\bf x$ with $8$-bit precision, ABSMAX quantization is defined as:
% \begin{align}
% \text{Quant}({\bf x}, 8) = \text{round}\left(\frac{{\bf x}}{\max(|{\bf x}|)} \cdot (2^{8-1} - 1)\right),
% \end{align}
% This relies on the max value within the input activation vector, so the find max operator is separated in Fig. \ref{Fig:Bitnet}.
% Afterward, Rotary Position Embedding (RoPE) is applied to query ($\bf Q$) and key ($\bf K$) vectors after the ternary linear layer. For a vector $\bf x$ at position $m$ and frequency $\theta_i = 10000^{-2i/d}$, the transformation for an element ${\bf x}_i$ in the first half is:
% \begin{equation}
% \begin{cases}
% x'i = x_i \cos(m\theta_i) - x{i+d/2} \sin(m\theta_i) & \text{for } i < d/2 \\
% x'{i+d/2} = x{i+d/2} \cos(m\theta_i) + x_i \sin(m\theta_i) & \text{for } i < d/2
% \end{cases},
% \end{equation}
% During the initial processing of a sequence of length $L$ (prefill stage), attention is calculated as:
% \begin{align}
% \text{Attention}({\bf Q}, {\bf K}, {\bf V}) = \text{softmax}\left(\frac{{\bf Q}{\bf K}^T}{\sqrt{d_h}}\right) {\bf V},
% \end{align}
% The computed $K$ and $V$ matrices are stored in the KV cache.
% For each subsequent token (decoding stage), the query $q_i$ uses the cached keys and values:
% \begin{align}
% \text{Attention}(q_i, K_{\text{cache}}, V_{\text{cache}}) = \text{softmax}\left(\frac{q_i K_{\text{cache}}^T}{\sqrt{d_k}}\right)V_{\text{cache}},
% \end{align}
% The computational load per decoding step is $O(Ld_{\text{model}})$, significantly lower than the prefill stage's $O(Ld^{2}_{\text{model}})$ due to the KV cache. New key and value vectors are appended to the cache after each step.
% Then, for a vector $x$ of dimension $d$, RMSNorm is calculated as:
% $\text{RMSNorm}(\bf{x}) = \frac{{\bf x}}{\sqrt{\frac{1}{d} \sum_{i=1}^{d} x_i^2 + \epsilon}} \cdot {\bf g}$,
% where $\bf{g}$ is a learnable gain parameter and $\epsilon$ is a small constant for numerical stability.
% The SwiGLU activation function in the subsequent FFN layer is defined as:
% $
% \text{SwiGLU}(\bf{x}, \bf{W}, \bf{V}) = \text{SILU}(\bf{x}\bf{W}) \odot (\bf{x}\bf{V})
% $,
% where $\odot$ denotes element-wise multiplication and $\text{SILU}(y) = y \cdot \sigma(y)$.
% Finally, a residual connection is applied around sub-layers of Attention/FFN:
% $
% {\bf Output} = {\bf x} + F({\bf x})
% $,
% where $\bf{x}$ is the input to the sub-layer before the RMSNorm and $F(\bf{x})$ is the function computed by the sub-layer.


% \subsection{Binary, Ternary, and Low-Bit Quantized Transformers}

% Model quantization is a key technique for compressing LLMs to run on constrained devices. Most conventional quantization approaches target 8-bit or 4-bit representations, but recent work has pushed the boundary down to the binary regime.

% \textbf{BitNet}~~\cite{bitnet} introduced a method for training Transformers from scratch using 1-bit weights. Despite extreme quantization, BitNet maintained competitive perplexity through custom scaling and layer-wise normalization strategies. Building on this, \textbf{BitNet-1.58}~~\cite{bitnet158} introduced ternary weight representations ($\{-1, 0, +1\}$), striking a balance between expressiveness and compression. Both approaches highlight the potential of binary LLMs in terms of storage, throughput, and energy efficiency. Similarly, \textbf{FBI-LLM}~~\cite{fbillm} and \textbf{OneBit}~~\cite{onebit} demonstrate fully binarized models trained using autoregressive distillation, achieving promising results on open-domain generation tasks. \textbf{DeepSeek-R1}~~\cite{deepseek} presents a hybrid quantization strategy applying ternary quantization to Mixture-of-Expert (MoE) layers, achieving up to $80\%$ model size reduction on a 671B model. Beyond training-time quantization, post-training quantization (PTQ) also plays a role. \textbf{BitDistiller}~~\cite{bitdistiller} combines self-distillation with quantization-aware techniques to push 3-bit and 2-bit LLM performance to new levels. \textbf{QuIP}~~\cite{quip} introduces 2-bit quantization with incoherence processing and rounding guarantees.

% These works focus on algorithmic aspects. In contrast, \textbf{TeLLMe} provides a hardware-aligned solution for deploying such models in practice, offering both matrix multiplication reuse and pipeline fusion for edge execution.

\subsection{Transformer/LLM Acceleration on Edge}
%\vspace{-1mm}
Deploying Transformers on FPGAs is challenging due to limited bandwidth and logic resources. Several works have explored quantized Transformer accelerators on embedded FPGAs/CPUs. 
\textbf{T-MAC} \cite{tmac} implements a TLMM kernel for CPUs using low-bit weights and high-bit activations. It achieves notable performance on edge CPUs. Moreover, the prefill speed is not ideal and close to the decoding throughput, suffering from the lack of parallelism computation bound inside edge CPUs. \textbf{{Li et al.}}~\cite{li2025pushing} successfully implemented a 4-bit quantized LLaMA2-7B model on the AMD KV260 platform. Although the model weights are quantized to 4-bit, the decoding computations rely on unquantized FP16 activations, thereby requiring all operations to be conducted in FP16. Moreover, hardware acceleration is limited to the decoding stage and does not address the computational demands of the prefill stage. \textbf{LLaMAF}~~\cite{LLaMAf} targets TinyLLaMA model with int8 quantization on ZCU102 with speed of 1.5 tokens/s. It leverages pipelined matrix-vector units on FPGA and asynchronous scheduling for non-linear operators on ARM A53 cores, but suffers from ARM-FPGA communication and ignore the prefill stage. \textbf{Edge-MoE}~\cite{edgemoe} introduces a memory-efficient MoE vision transformer accelerator using dynamic task-level sparsity. 
A key technique is the specialized reordering to enable the data reuse of attention computation for edge FPGAs when the attention score matrix is dense. \textbf{SECDA}~\cite{haris2024designing} designs the matrix multiplication accelerator supporting block FP quantized operations on PYNQ, reducing latency by 11x compared to dual-core Arm NEON-based CPU execution for the TinyLLaMA model. However, the tokens/s is only 0.58, which means 2 seconds for one token generation.
\textbf{MEADOW}~\cite{moitra2025meadow} presents a memory-efficient weight packing framework for low-power edge LLMs, implemented on the Xilinx ZCU102 FPGA, with approximately 10W power consumption. It introduces a weight packing strategy for memory-efficient linear operation in both prefill and decoding stages. However, they only implement the INT8 matmul without considering the end-to-end acceleration of FP16 smooth and dequant operation that supports the utilized INT8 smoothquant~\cite{xiao2023smoothquant}.  
It reaches the prefill speed up to approximately 0.64 ms at 64 tokens and decoding speed up to 2 token/s for the OPT 1.3B model from the profiling result. Compared to these, our work is the first to unify binary weight inference, high throughput prefill/decoding support, and FPGA memory hierarchy optimization into one cohesive design.
% \todo[inline]{
% \textbf{\textcolor{blue}{MEADOW}} \url{https://arxiv.org/pdf/2503.11663v1}

% \textbf{\textcolor{blue}{LightMamba}} \url{https://arxiv.org/abs/2502.15260}
% \textbf{\textcolor{blue}{TerEffic}}
% \textbf{\textcolor{blue}{SlimLLaMA}}
% }

% \textbf{\todo[inline]{Edge-LLM}}~~\cite{edgellm} proposes a heterogeneous CPU-FPGA acceleration framework for efficient LLM inference on edge platforms. The system integrates several optimizations, including structured sparsity, asynchronous data transfers, and group vector systolic arrays to enhance hardware utilization. It supports hybrid precision computations (e.g., FP16 × INT4 and FP16 × FP16) through customized hardware operators. An end-to-end dynamic compilation scheme is introduced to map diverse LLM workloads onto the heterogeneous architecture.



\subsection{LLMs on High-performance FPGAs}
%\vspace{-1mm}
\textbf{LightMamba}~\cite{LightMamba} introduces an FPGA-friendly post-training quantization algorithm for the Mamba2-2.7B model, supporting W4A4 quantization. Deployed on high-end Versal VCK190 and U280 FPGAs, it achieves throughputs of 7.21 and 93 tokens/s, respectively.
\textbf{FlightLLM}~\cite{FlightLLM} deploys the LLaMA2-7B model on Alveo U280 and Versal VHK158 FPGAs, outperforming NVIDIA GPUs in single-batch inference with a throughput of 153 tokens/s.
\textbf{EdgeLLM}~\cite{EdgeLLM} utilizes the VCU128 FPGA with HBM to deploy LLaMA2-7B, achieving a higher decoding speed than FlightLLM.
\textbf{TerEffic}~\cite{yin2025tereffic} deploys a ternary 1.3B non-transformer model on the Alveo U280, achieving a remarkable 1400 tokens/s. However, its reliance on soft logic for ternary matrix multiplication, rather than a table-lookup approach, results in over 700K LUTs consumed for redundant paralleled ternary matmul.
These works leverage high-end FPGAs with over 500K LUTs, thousands of DSP units, thousands of BRAM cells, HBM providing 460 GB/s bandwidth, and power consumption exceeding 50W, with costs exceeding \$10,000 (vs. \$248 for the KV260)~\cite{AMDVersalOverview}. While these platforms maximize FPGA performance, their performance is not directly comparable to edge FPGAs.
%\vspace{-1mm}
\subsection{Existing Table Lookup Works on FPGAs}
%\vspace{-1mm}
Using DSP48E1 blocks for 1.58-bit TLMM on FPGAs is highly inefficient due to their design for high-throughput, fixed-point arithmetic (25x18-bit multiplier, 48-bit accumulator). Ternary operations, with weights of -1, 0, or 1, require only simple selection or negation, underutilizing the DSP’s complex logic. Instead, LUTs can efficiently handle “pass,” “negate,” or “zero” operations, making DSPs a resource- and power-inefficient choice for highly quantized neural networks better suited to FPGA fabric logic.
At a larger architectural scale, implementing ternary matmul with purely selection-based logic also becomes inefficient in terms of resource utilization and latency. While ternary weights simplify multiplication into selections, achieving high parallelism demands numerous LUT-based multiplexers. This approach is suboptimal for large models (i.e., with large $d_{\text{model}}$), where the limited set of ternary operations results in repetitive computational patterns \cite{TLMAC}. 
% For instance, with three activation values, there are only 27 possible combinations to store in the TL table. In a model with $d_{\text{model}}=2048$, many of these combinations recur in the TL table, leading to redundant computations. This repetition reduces the efficiency of a straightforward lookup-based approach, as the same operations are unnecessarily repeated.
As an alternative to logic-based acceleration, a table-lookup method has been proposed as an efficient solution for ternary matmul on CPUs ~\cite{tmac}. This technique leverages specialized SIMD instructions, such as ARM NEON and AVX, for rapid parallel look-ups. However, its effectiveness is often constrained by the limited size of the look-up table (256 bit/128 bit), which can lead to frequent, high-latency memory accesses when the table is swapped.

FPGAs provide an optimal solution by leveraging their inherent LUT resources to implement larger tables using LUT-based Distributed RAM. As demonstrated in studies such as~\cite{LUTNET, LOGICNET,SUMLUTNET}, each basic LUT can be configured to store values that facilitate neural network inference. However, fine-grained control of LUT resources is currently limited to simple tasks like MNIST and smaller network sizes in the thousands level of neurons. To expand the area of TL-based computing on FPGA,~\cite{TLMAC} introduces a CNN-specific TL engine that adopts a coarse-grained approach to distributed LUTRAM. Yet, this approach focuses solely on TL convolution logic, overlooking comprehensive dataflow considerations for real-world INT/FP mixed-precision quantized neural network workloads. Moreover, they lack the design in TL logics' interactions with on-chip and off-chip smemory, which is vital in the memory bound trait of transformer and LLM inference. To the domain to the LLM era, we propose a TLMM engine for LLM inference on FPGAs, integrated with a detailed exploration of efficient mix-precision design and fine-grained on-chip URAM management. 
% \subsection{Importance of Prefill Latency} \todo{write with discussion and examples}

\begin{figure}[t] 
\centering
\includegraphics[width=\linewidth]{Figures/TeLLMe_FIG_1_new.pdf}
\caption{Breakdown of TeLLMe 1.58-bit Model Inference Process with Prefill and Generation}
\label{Fig:Bitnet}
\end{figure}
% \todo[inline]{
% 1. The Bitnet module is not correct
% 2. We should add the tag for which units are fused together
% }
% % \section{BitNet 1.58B Architecture}
% \subsection{Background of Ternary matrix multiplication}
% \subsection{Background of PD Disaggregation}

\section{TeLLMe Design Methodology}
%\vspace{-1mm}
% \subsection{TeLLMe Overview}
% \todo[inline]{Talk about the overall challenge and solution in this project
% Goal: Achieve an efficient Tenantry Language Model Inference on a low-end edge FPGA
% Challenge: 
% \begin{itemize}
%  \item Inefficient DSP mapping for W1.58A8 matrix multiplication 
%  \item Hybrid Precision Support for Quantization
%  \item Intrinsic different computation pattern of decoding attention and prefill attention on FPGA 
% \end{itemize}
% Solution:
% \begin{itemize}
%     \item Efficient TLMM design  
%     \item Fine-grained URAM weight buffer management
%     \item Streaming fusion of elementwise operation and TLMM
%     \item Choice of Operator Fusion for supporting ABSMAX quantization
%     \item Memory Efficient reversed-reordered prefill attention and On-chip memory supported decoding attention
% \end{itemize}
% }
% The primary objective of this project is to implement an efficient framework for ternary large language model (LLM) inference on resource-constrained low-end FPGA platforms. The main challenges are presented as follows,

% As depicted in the Fig. \ref{fig:sys_arch}, we mainly have below modules for the accelerator design: (1) Look-up table-based ternary matrix multiplication for both decoding and prefill pass (2) Specialized Reverse Reorder for Prefill attention (3) Unified decoding attention and LM-Head (4) Special function units.  
% As shown in Fig. \ref{fig:sys_arch}, the accelerator design primarily consists of the following modules: (1) a table-look-up-based ternary matrix multiplication for both the decoding and prefill passes; (2) a specialized reverse reorder for prefill attention; (3) a unified decoding attention and language model head (LM Head); and (4) specialized functional units.
% \todo[inline]{List of used symbol}

% \begin{table}[t]
% \centering
% \caption{Key Parameters and Their Descriptions in Our Design}
% \label{tab:parameters}
% \begin{tabular}{lp{6cm}} 
% \toprule
% \textbf{Symbol} & \textbf{Description} \\
% \midrule
% $d_{\text{model}}$ & The dimensionality of the model's hidden states. \\
% $d_{\text{head}}$ & The dimensionality of a single attention head. \\
% $d_{\text{ffn}}$ & The dimensionality of \\ &the feed-forward network's intermediate layer. \\
% $q,k,v,o$ & Identifies the query, key, value and post-attention output projection\\
% $m, n, k$ & Dimensions for a general matrix multiplication (${\bf C}_{m \times k} = {\bf A}_{m \times n} \cdot {\bf B_{n \times k}$). \\
% ${\bf W}_{\text{idx}}$ & A weight matrix containing multiple packed ternary weights. \\

% $G$ & The number of ternary weights packed into a single $W_{\text{idx}}$ word index. \\
% $B_{\text{idx}}$ & The bitwidth of a packed weight word index, $W_{\text{idx}}$. \\
% $B_{\text{TB}}$ & The bitwidth of a TL table entry. \\
% $N_{\text{TB}}$ & The number of entries in a single TL table. \\
% $S_{\text{TB}}$ & The size of all TL tables in bit. \\
% $T$ & The total number of concurrent TL tables instantiated for the lookup operation. \\
% $Q$ & The number of parallel lookups performed on a single TL table. \\
% $X_{\text{URAM}} \times Y_{\text{URAM}}$ & The dimensions of the URAM buffered weight matrix. \\
% $N_{\text{URAM}}$ & The total number of URAM blocks available on the target FPGA. \\
% $c_{\text{URAM}}$ & The number of URAM blocks cascaded together to support a  larger bitwidth of $T \times B_{\text{idx}}$\\
% \bottomrule
% \end{tabular}
% \end{table}

\begin{figure}[t] 
\centering
\includegraphics[width=\linewidth]{Figures/HLS_SYSARCH_YE.pdf}
\caption{System Architecture of TeLLMe}
\label{fig:sys_arch} 
\end{figure}

\subsection{Overview}
The technical challenges of ternary LLM implementation on edge FPGA are shown as follows, 
\textbf{\ding{182} Challenge of implementing TLMM for Ternary LLM on FPGA:} The TLMM approach is well-suited for implementing ternary linear operations. However, the current TL-based method on FPGA~\cite{TLMAC, LOGICNET} focuses solely on TL-based convolution or small-sized perception machine, overlooking its interaction with on-chip and off-chip memory. This interaction is critical for LLMs, which are memory-bound and heavily relying on efficient memory management.
\textbf{\ding{183} Hardware Support for Mixed-Precision in Low-bit Quantization:} Modern quantization strategies often employ a hybrid-precision approach, where weights and activations are represented with both FP16 (attention, dequant, quant, rotary embedding, etc) and quantized low-bit integers (linear). 
Designing a unified hardware architecture that can process these varied data types without incurring significant performance or area overhead is a complex task.
\textbf{\ding{184} Diverse Computational Patterns in Attention Mechanisms:} The attention mechanism in transformer models exhibits two different computational profiles. The initial \textbf{prefill} phase is compute-bound and characterized by parallel processing of the input prompt, while the subsequent \textbf{decoding} phase is memory-bound and involves sequential, auto-regressive token generation. A successful FPGA accelerator should efficiently handle both distinct patterns.

% To address these challenges, we present the following solutions:
% \ding{182} Efficient addition-only TLMM engine 
% \ding{183} Fine-grained URAM weight buffer management and DDR access management
% \ding{184} Streaming fusion of FP16-based element-wise operation for latency overlapping
% \ding{185} Optimal choice of operator fusion for supporting ABSMAX quantization
% \ding{186} Memory-efficient reversed-reordered prefill attention and on-chip memory-supported decoding attention

The overall Bitnet model architecture is depicted in Fig. \ref{Fig:Bitnet}, comprising attention computation, linear projections with quantization and dequantization, rotary position embedding (RoPE), RMSNorm, absolute maximum (ABSMAX) quantization, and SwiGLU modules. During the prefill phase, the KV projections are stored in the KV cache, whereas in the decoding phase, the cached KV values are loaded for attention computation.
As illustrated in Fig. \ref{fig:sys_arch}, the accelerator architecture comprises the following key modules to address the challenges: (1) a TLMM engine that handles both decoding and prefill passes, integrated with element-wise operations and quantization, referred to as the TLMM-FUSE unit to address \ding{182}, \ding{183}; (2) an on-chip weight buffer management unit (WBMU) that handles the off-chip DDR weight data transfer and accessing of the URAM weight buffer to address \ding{182}; (3) a reversed-reordered prefill attention (RPA) unit that get rid of the redundant attention mask during prefill stage to address \ding{183}, \ding{184}; (4) a decoding attention (DA) unit that match the vector-wise computation in decoding stage to address \ding{183}, \ding{184}; and (5) an RMSNorm unit combined with a find-max operation for supporting the ABSMAX quantization, refered to as RMS-MAX unit to address \ding{183}. 

All modules are designed in a dataflow style to enable function-level pipelining. The DDR address ports are bound to distinct MAXI bundles, facilitating independent data streaming from DDR memory. Four temporary address ports are allocated for intermediate storage in DDR. The weight address ports handle weight loading, while the K cache and V cache address ports manage the prefill attention KV return values and decoding attention input values, respectively. The relationship between these modules and the model's software logic is illustrated in Fig.~\ref{Fig:Bitnet}.


% \begin{figure}[h] 
% \centering
% \includegraphics[width=\linewidth]{Figures/lutmatmul_YE.drawio.pdf}
% \caption{High Level Dataflow of TLMM ($G = 3, T=3,Q=3$)}
% \label{fig:lutmatrix multiplication} 
% \end{figure}
% \subsubsection{TLMM Engine Design}\label{SEC:TLMM Design}

% \todo[inline]{SHOULD MAKE THE FIGURE BIGGER AND DESCRIBLE DETAIL ARCHITECTURE IN EACH SUB=UNIT}



% \section{TLMM, Floating Point Operation, and Weight Management Design}
% %\vspace{-1mm}

% \todo{should be moved to background}Look-up table based matrix multiplication has been one of the most efficient method in CPU design of ternary matrix multiplication, because of the provided special assembly of Look up table instruction in ARM NEON and AVX which supports 128/256bit look up table. However, the limited table size may requires more access of the look up table space for extra latency ~\cite{tmac}. To solve this problem, FPGA can be utilzied as the ideal component for larger size of look up table because of the intrinsic LUT-6 components. Only a few FPGA works explore this topic  ~\cite{TLMAC} and most of them are exploring on the module itself instead of any potential efficient dataflow and scheduling. Our work provides the Loop up table ternary multiplication free matrix multiplication implementation and efficient scheduling and pipeline exploration.

% \todo[inline]{Mention the FUSE module and WBMU design as well, and showcase the relevant Fig. }
% In this section, the overall units for the solution \ding{182}-\ding{183} to support efficient TLMM operation of the edge FPGA are introduced. The overall TLMM-FUSE system, illustrated in Fig. \ref{fig:TLMM_FUSE}, is built upon a streaming dataflow paradigm to maximize throughput. To decouple the processing modules and ensure smooth data transfer, all internal communication channels utilize stream FIFOs. The system interfaces with off-chip DDR memory, where a dedicated read stream unit fetches and converts memory-mapped data into streams for on-chip computation and writes back through the write stream unit. The WBMU is responsible for managing the retrieval and distribution of model weights.


\subsection{Table-Lookup Matmul (TLMM) Engine}
%\vspace{-1mm}
% \todo[inline]{Remember to unify the symbol between figure and text}
% \todo[inline]{We should merge Fig. 3b, 3c to Fig. 4}


% \todo[inline]{Add the description in a more detailed way to explain why generic DSP-based matrix multiplication on FPGA is not efficient }


% \todo[inline]{Include the padding logic stemming from the OBFU unit}
% As described in Fig. \ref{fig:lutmatrix multiplication}, the TLMM can be divided into two stages: (1) preprocessing the weights into groups sized $G$, and (2) performing the online ternary matrix multiplication computation. 
% In the preprocessing stage, assume that every $G$ ternary values are packed into a single index for TL table addressing, resulting in \(N_{\text{TB}} = 3^{G}\) combinations. The index representation for this packing requires \( B_{\text{idx}} = \lceil \log_2 N_{\text{TB}} \rceil \) bits. Let \( {\bf A} \in \mathbb{N}^{m \times n} \) and \({\bf W} \in \text{ternary}^{n\times k} \). The preprocessing of the weights involves encoding every group of \( G  \) ternary values into a $B_{\text{idx}}$-bit packed index.
% In the online stage, we perform vector-wise tiled matrix multiplication between \( {\bf A} \) and \( \bf W \). The first step is to fetch a FIFO data of $T \times G$ INT8 activation and establish the LUT distributed RAM supported TL table using the precompute unit, which consists of a set of adders and subtractor trees to cover all possible $N_{\text{TB}}$ combinations of dynamic activations. Each TL table entry bitwidth is sized $B_{\text{TB}} = 8 + log_{2}\lceil G \rceil$ as this is the bitwidth of the maximum value of $G$ INT8 data reduction. The packed-up word indices are then used to address and select the corresponding values from the TL table. Finally, the selected values are accumulated to generate a single vector \( {\bf O} \in \mathbb{N}^{1 \times k} \) in the output matrix and are pushed into the FIFO channel for later FP16 operations. The Fig. \ref{fig:lutmatrix multiplication} showcases when $G = 4$, $T =4$, $Q = 3$ the situation. The hardware detailed design and detailed input and output connection is dipicted in Fig. \ref{fig:TLMM_FUSE} . 

\subsubsection{TLMM Engine Design} \label{sec:TLMM design}
The TLMM method, illustrated in Fig. \ref{fig:lutmatrix multiplication}, consists of two main stages: offline weight preprocessing and online matrix multiplication. In the offline stage, the ternary weight matrix \({\bf W} \in \{-1, 0, 1\}^{n \times k}\) is partitioned into groups of size $G$, which are then encoded into compact indices. This encoding yields \(N_{\text{TB}} = 3^{G}\) unique combinations per group, requiring an index bitwidth of \( B_{\text{idx}} = \lceil \log_2(3^G) \rceil \) bits.

During the online stage, the matrix-vector product between the activation matrix \({\bf A} \in \text{INT8}^{m \times n}\) and the preprocessed weight indices is computed. For each output vector $\bf o$, a group of $G$ INT8 activation values is fetched to dynamically generate a small TL table implemented by distributed RAM. This table is populated by a pre-computation unit of adder/subtractor trees and stores all \(N_{\text{TB}}\) possible partial sums for the current activations, with each entry sized at a bitwidth of \(B_{\text{TB}} = 8 + \lceil \log_2 G \rceil\) to prevent overflow. The preprocessed weight indices are then used to look up these partial sums from the TL table. To support the parallel multiple read requests of $Q$ of $T$ TL tables, each TL table has to duplicate $Q$ times, resulting in the TL table size of $B_{\text{TB}} \times N_{\text{TB}}\times Q \times T$. Finally, these values are accumulated to produce the output vector \({\bf o} \in \text{INT32}^{1 \times k}\), which is pushed into the stream first-in-first-out unit (FIFO) channel for subsequent operations. 
To better vectorize the TLMM and facilitate efficient on-chip URAM access, the consecutive \( T \) indices can be grouped into an index vector, enabling simultaneous access to different TL tables. The weight index vector matrix can be rewritten as \( {\bf W}_{\text{idx}} \in \{B_{\text{idx}}, T\}^{\frac{n}{T\times G} \times k} \). 


\begin{figure}[h] 
\centering
\includegraphics[width=\linewidth]{Figures/lutmatmul_YE.drawio.pdf}
\caption{High Level Dataflow of TLMM ($G = 3, T=3,Q=3$)}
\label{fig:lutmatrix multiplication} 
\end{figure}
Regarding the scheduling in Fig. \ref{fig:lutmatrix multiplication} with $G=4$, $T=3$, and $Q=3$ and the detailed architecture in Fig. \ref{fig:TLMM_FUSE}, the innermost loop first performs vector operations to establish the \( T \) TL tables through the precompute tree in parallel, based on the first unpacked \( T \times G \) entries of the \(\bf A \) matrix from the stream FIFO channel. Then, leveraging the multiple reading traits of multiple independent partitioned URAM blocks (array partition strategy and parameter will be detailed in the \ref{SEC:WBMU}), \( Q \) \( {\bf w}_{\text{idx}} \in {\bf W}_{\text{idx}} \) index vectors are processed in parallel for TL table addressing, returning \( Q \times T \) outcomes. The corresponding TL table return values are then accumulated in INT32 into an output buffer of size \( k \). The \( k \) index vectors on each row of \( W \) are traversed in steps of \( Q \). The TL table addressing and accumulation process can be fully pipelined with an interval of one cycle, as there are no inter-iteration dependencies. After the first \( T \times G \) entries of the \( \bf A \) matrix are processed, the \( m \) values of the row of \( \bf A \) are traversed in steps of \( T \times G \) in the intermediate loop. Finally, the outermost loop traverses the different channels in \(\bf A \), corresponding to the tokens in the prefill stage of the LLM. 

%The detailed hardware design is depicted in Fig. \ref{fig:TLMM_FUSE}. The input data $T\times G$ INT8 activation is first unpacked into the precomputation tree of addition or subtraction to construct the $T$ TL tables.
% Then the 

 As for the interaction with the weight buffer, \( {\bf W}_{\text{idx}}\) is loaded onto the on-chip URAM in a single pack of DDR loading requests for each ternary linear operation and fully decouple with the computation of TLMM to prevent the compute engine from stalling. Moreover, our TLMM engine supports 3 sizes in the TLMM: $q,k,v,o$ projection sized $d_{\text{model}} \times d_{\text{model}}$, up and gate projection in FFN sized $d_{\text{model}} \times d_{\text{ffn}}$, and down projection in FFN sized $d_{\text{ffn}} \times d_{\text{model}}$. This is also supported by the WBMU to convert the high-level address in weight matrices into a physical URAM access address, and additional zero-padding weight indices will be added to support the $T \times G$ parallelism and transpose relationship between up and down projection. The detailed architecture will be presented in Sec. \ref{SEC:WBMU}.

\subsubsection{Comparison between Different Ternary Matrix Multiplication Designs} \label{sec:TLMM_METHODS}
% In general, there are three ways to implement the ternary matrix multiplication using pure fabric block. The first one is the naive implemenation straightfoward select logic of whether keep pass or negative the relevant value. The second one is partial storing the value in the TL table as the symmetric trait that the -1 and 1 of ternary weight and than the $N_{\text{TB}} = (3^{G} - 1)/2$, where the minused one is the all zero state. However the implment this, we also have to have additional all unroll resource of identifiing whehter the index is negative and revert them to surpport $T \times Q$ parrelell read. The third way is our implementation to completely store all the possible entry with $N_{\text{TB}} = 3^{G} $, in this way, we don't need to instantiate the additional logic and pure look up addressing for $T \times Q$ parrelell read. For TL table itsself, the depth have to meet the distributed RAM mapping tech \todo[inline]{cite user guide}, for example the basic unit is RAM32X1 on zynq ultrascale+ device, as minuim depth is 32 and the cascade will be a multiple of $2^5$. In our design the group is configured at $G = 3$ the minimize the resouce on the edge FPGA, $(3^{3} - 1)/2 < 3^{3} < \times 2^{5}$, the LUT for the two TL tables will be the same. In the same time, our implemention save the resource in parrellel read thus more effcient. 
% Then the lut consuption logic can be charactized. We assume the LUT resouce for parrell read supporting logic is $\text{LUT}_{p}$, the resource for the generation of TL table is $\text{LUT}_{gt}$, the resource for TL table logic cell its self is $\text{LUT}_{t}$. 


% In general, there are three methods for implementing TLMM using the FPGA fabric logic.
% The first is a naive implementation that utilizes straightforward selection logic to pass, negate, or zero-out the relevant input value according to the ternary weight.
% The second method involves partially storing results in a TL table. This approach leverages the symmetric property of ternary weights (+1, -1), storing entries only for the positive values. The required number of entries in the TL table is $N_{\text{TB}} = (3^{G} - 1)/2$, where $G$ is the group size and the subtraction of one accounts for the all-zero state. However, to support $T \times Q$ parallel reads, this implementation requires additional logic to identify negative indices and invert the corresponding outputs from the table.
% The third method is our implementation, which completely stores all possible entries in the TL table, for a total of $N_{\text{TB}} = 3^{G}$. This way, we eliminate the need for additional selection logic and can rely on a pure look-up addressing scheme for $T \times Q$ parallel reads. 
% The LUT utilized for precompute can be estimated as $LUT_{\text{pre}} = T \times N_{\text{TB}} \times LUT_{\text{tree}}$. The size of TL table is formultated as $LUT_{\text{TL}} = Q\times T  \times N_{\text{TB}} \times LUT_{\text{entry}}$, where $LUT_{\text{entry}}$ is LUT utilze to store one entry. The size of lookup logic is formulated as $LUT_{\text{LPL}} = Q\times T \times LUT_{\text{LP}}$, where $ LUT_{\text{LP}}$ is LUT consumed in one look up and corresponding conversion logic. 
% In our case of highly parallel operation of $T \times Q$ in method 2,3 look-up or $T \times G \times Q$ selection in method 1, the $LUT_{\text{LPL}}$ is the conducting factor as the logic of storing is intrinic for $LUT_{TL}$ architecture of distributed logic of RAM32X1 ~\cite{AMD_UG974_ultrascale}. The final result is supported by\ref{SEC:ABLATION_TLMM}. 
In general, three methods can be employed to implement TLMM using FPGA fabric.
Method 1 is a naive implementation that employs straightforward selection logic to pass, negate, or zero out the input value based on the ternary weight.
Method 2 partially stores results in a TL table. By exploiting the symmetry of ternary weights (+1, –1), only entries corresponding to positive weights are stored. The number of required entries is given by \( N_{\text{TB}} = (3^G - 1)/2 \), where \( G \) is the group size and the subtracted unit accounts for the all-zero state. However, to support \( T \times Q \) parallel reads, this scheme requires extra logic to identify negative indices and invert the corresponding outputs retrieved from the table.
The third method—our implementation—stores all possible entries in the TL table, totaling \( N_{\text{TB}} = 3^G \). This eliminates the need for additional selection logic and enables a purely lookup-based addressing scheme for \( T \times Q \) parallel reads.
As for resource analysis, the LUT resources for precomputation can be estimated as  
\begin{align}
\text{LUT}_{\text{PRE}} = T \times N_{\text{TB}} \times \text{LUT}_{\text{tree}},
\end{align}
where the $\text{LUT}_{\text{tree}}$ is the average LUT utilized in each output of the precomputation tree.
The size of the TL table is formulated as  
\begin{align}
\text{LUT}_{\text{TB}} = T \times Q \times N_{\text{TB}} \times \text{LUT}_{\text{entry}},  
\end{align}
where \( \text{LUT}_{\text{entry}} \) is the number of LUTs needed to store one entry in TL table. The lookup logic LUT consumption is given by  
\begin{align}
\text{LUT}_{\text{LPL}} = T \times Q \times \text{LUT}_{\text{lp}},
\end{align} 
where \( \text{LUT}_{\text{lp}} \) represents the LUT cost of a single lookup, its associated conversion logic, and reduction logic.
In the highly parallel \( T \times Q \) lookups in Methods 2, the dominant factor is \( \text{LUT}_{\text{LPL}} \), since the $N_{\text{TB}} \times \text{LUT}_{\text{entry}}$ is efficiently implemented using the distributed RAM architecture of primitives such as cascaded RAM32X1~\cite{AMD_UG974_ultrascale}. These argument are further validated in the ablation study  Sec.~\ref{SEC:ABLATION_TLMM}.

% Specifically, the naive and partial-storage methods necessitate supplementary selection and inversion logic, leading to greater LUT utilization, as will be demonstrated in the ablation study in Sec. \ref{SEC:ABLATION_TLMM}.
% Regarding the hardware implementation of the TL table itself, its depth must conform to the distributed RAM mapping capabilities of the target device. For a Zynq Ultrascale+ device, the basic unit is a RAM32X1, which has a depth of 32 ~\cite{AMD_UG974_ultrascale}. Larger memories are cascaded from this primitive, making their effective depth a multiple of 32 ($2^5$).
% In our design, we configure the group size at $G=3$ to constrain resource usage on edge FPGAs. For the partial storage method, this requires $(3^3 - 1)/2 = 13$ entries, while our full storage method requires $3^3 = 27$ entries. Thus, the real number of entries of TL table is supposed to be $N'_{\text{TB}} = 32\lceil N_{\text{TB}}/32\rceil$. Since both `13` and `27` are less than the minimum primitive depth of `32`, the number of LUTs consumed by the TL table itself is identical for both approaches. At the same time, our implementation saves the resources otherwise needed for the parallel read logic, making it more efficient.



% \begin{figure}[t] 
% \centering
% \includegraphics[width=\linewidth]{Figures/TLMM_FUSE_UNIT_YE.pdf}
% \caption{Detailed design of TLMM-FUSE Unit. \todo[inline]{INPUT OUTPUT WRONG}}
% \label{fig:TLMM_FUSE} 
% \end{figure}



% \begin{figure}[t] 
% \centering
% \includegraphics[width=\linewidth]{Figures/Rope.pdf}
% \caption{Detailed design of TLMM-FUSE Unit. \todo[inline]{INPUT OUTPUT WRONG}}
% \label{fig:Rope} 
% \end{figure}

% \begin{figure}[t] 
% \centering
% \includegraphics[width=\linewidth]{Figures/TLMM_schedule.pdf}
% \caption{Dataflow overlapping of TLMM-FUSE}
% \label{fig:TLMM_schedule} 
% \end{figure}
\captionsetup{font=small}
\captionsetup[subfigure]{font=footnotesize,labelformat=parens,labelsep=space}
\begin{figure}[t]
    \centering
    \begin{subfigure}[b]{0.47\textwidth}
        \centering
        \includegraphics[width=\linewidth]{Figures/TLMM_FUSE_UNIT_YE.pdf}
        \vspace{-2mm}
        \caption{Detailed Design of TLMM-FUSE Unit}
        \label{fig:TLMM_FUSE}
    \end{subfigure}
    \hfill % Adds horizontal space between subfigures
    \begin{subfigure}[b]{0.27\textwidth}
        \centering
        \includegraphics[width=\linewidth]{Figures/Rope.pdf}
        \vspace{-2mm}
        \caption{Mem. Dependency of Different RoPE} % Corrected caption for clarity
        \label{fig:Rope}
    \end{subfigure}
    % \hfill % Adds horizontal space between subfigures
    \begin{subfigure}[b]{0.20\textwidth}
        \centering
        \includegraphics[width=\linewidth]{Figures/TLMM_schedule.pdf}
        \vspace{-2mm}
        \caption{Dataflow Overlapping}
        \label{fig:TLMM_schedule}
    \end{subfigure}
    \vspace{-2mm}
    \caption{Overview of TLMM hardware components and dataflow. (a) The TLMM-FUSE unit design. (b) The comparison of different RoPE operations. (c) The dataflow execution schedule.}
    \label{fig:combined_designs}
\end{figure}

\subsection{Fused Element-wise Operations}
%\vspace{-1mm}
% Elementwise operation denotes the operation that requires one or two sets of input operated in the elementwise manner, every element is consumed only once with no inter-depnedency. Inside Bitnet module,:  (1) Quantazation from FP16 value to INT8 value   (2) Dequantization from INT32 value to FP16 value  (2) Rope operator the addition between rope cache and computed QK value. (3) Elementwise Add for residual connection (4) Elementwise multiplication (5) SILU operation before it. Previous work seldom care about these operation and regard them as negligible and none resource consuming.  (7) Find column-wise max in ABSMAX quantzation. However, in real implmentation, these operator somehow bound the overall latency and resourse utilization. In this paper we carefully design the fusion unit to overlap their latency with the linear operation.
% Element-wise operations refer to computations performed on one or two input tensors, where each element is processed independently and consumed only once. Within the Bitnet module, several categories of such operations are critical: (1) quantization from FP16 to INT8;  (2) the column-wise maximum value looking up in ABSMAX quantization. (3) dequantization from INT32 to FP16; (4) the Rotary Position Embedding (RoPE) operator, which involves pair-wise multiplication and addition between the RoPE cache and the computed Q and k values; (5) element-wise addition for residual connections; (6) element-wise multiplication; (7) the preceding activation function.

% Previous research has largely overlooked the performance impact of these operations, often considering them negligible and not resource-intensive. In practical deployments, however, these operators can become a performance bottleneck, significantly constraining the model's overall latency and resource utilization, especially within a quantized module which invlove both efficient INT only operation and demanding FP operations.  Moreover, unlike CPU and GPU operation, the ALUs can be reused for these operations only latency increased; every mult and add for these operation refers to the modules own DSP and LUT computing resource, leading to more resource consumption. To address this issue, this paper introduces a meticulously designed fusion unit aimed at overlapping the latency of these element-wise computations with the core linear operations, thereby optimizing execution efficiency.
Element-wise operations are computations performed on input tensors where each entry or entry pair is consumed only once. Within the Bitnet architecture, several such operations are critical: (1) quant from FP16 to INT8 and subsequent dequant from INT32 to FP16; (2) the channel-wise maximums for ABSMAX quantization; (3) the RoPE operator; (4) element-wise additions for residual connections; (5) element-wise multiplications in SwiGLU; and (6) the activation functions.
Previous research on LLM accelerators has overlooked the performance impact of these operations, often deeming them negligible \cite{moitra2025meadow}. In practical deployments, however, these operators can create performance bottlenecks, constraining the model's overall latency and resource utilization. This issue is acute in quantized modules, which involve a mix of efficient high-throughput INT computations and demanding FP operations. 
% Furthermore, the implementation on FPGAs differs fundamentally from that on CPUs or GPUs. While CPUs and GPUs reuse general-purpose Arithmetic Logic Units (ALUs), FPGA-based designs require dedicated DSP and LUT resources for each operation, leading to greater resource consumption.
To address the above challenges, we introduce a specialized fusion unit designed to overlap the latency of these element-wise computations with the core TLMM operations, thereby optimizing execution efficiency. Moreover, a vector-wise operation is implemented to make sure we match the throughput of TLMM under resource constraints of DSP and LUT when designing FP operators. The detailed streaming pattern and datapath are displayed in Fig. \ref{fig:TLMM_FUSE}. The stream FIFO channels are in grey color with the type of vector it is buffering.
% \todo[inline]{Add the Fig. description of the arbiter stream path and vectorize/fifo resizing}

\subsubsection{INT8 Quantization and FP16 Dequantization}
% To seamlessly fuse the quantization with linear operation, the stream fifo channel is connected between these unit. Moreover, all the FP16 data are all packed up into the 256-bit FP16 vector to enable the paralleled computing and efficient DDR access of 256bit AXI DATA. As for quantazation, the packed vector size is changing from 16 FP16 data to $T \times G$ INT8 data for TLMM input. Thus, the stream resizing logic is involved in quantization. For dequantazaion, the values are undergoing 16 FP16 vectorized reciprocal and multiplication.
To seamlessly fuse quantization with linear operations, a stream FIFO channel connects the respective functional units. To maximize data throughput, all FP16 values are packed into 256-bit vectors, each holding 16 elements. This vectorization enables parallel computation and efficient 256-bit wide AXI access to DDR memory. The quantization process involves a data-width transformation, converting these 16-element FP16 vectors into $T \times G$ INT8 vectors for the TLMM input FIFO, which necessitates stream resizing logic. Conversely, dequantization performs vectorized multiplication operations on the 16-element INT32 vectors to convert them to FP16 precision.

% \todo[inline]{Add the Fig. description of the arbiter stream path and vectorize/fifo resizing}

\subsubsection{Elementwise Activation-fused Multiplication and Addition}
The element-wise multiplication with SILU, element-wise addition, and RoPE operations are all performed in floating-point precision following dequantization. Both these three operations and bypass logic are implemented in TLMM-FUSE as displayed in Fig. \ref{fig:TLMM_FUSE} with an arbitrary unit deciding which submodule the stream FIFO flow in. As shown in Fig. \ref{Fig:Bitnet}, the element-wise multiplication is an integral component of the SwiGLU. This function operates on the outputs of the gate and up projections: the gate projection's output is first processed by a SILU activation, and the result is then multiplied element-wise with the up projection's output. The element-wise addition is applied immediately before the RMSNorm layer to incorporate the residual connection.

% For consideration of LUT resource efficiency, the SILU activation is implemented by the accurate multi-segment approximation introduced in \todo[inline]{cite}. 


\subsubsection{RoPE Operation} 
For the RoPE implementation, we assume the necessary sinusoidal values are pre-computed and stored in DDR, as generating them on-the-fly would be resource- and time-intensive. The core RoPE function operates on a vector \( \mathbf{x} \in \text{FP16}^{d_h} \) by applying a rotation to pairs of elements. The mathematical formulation of the two implementations differs only in how the vector indices are grouped for these rotations:
\begin{equation}
\begin{cases}
\tilde{x}_m^{(2t)} = x_m^{(t)} \cos(m\theta_t) - x_m^{(t + d_h/2)} \sin(m\theta_t) \\
\tilde{x}_m^{(2t+1)} = x_m^{(t + d_h/2)} \cos(m\theta_t) + x_m^{(t)} \sin(m\theta_t)

\end{cases}\label{eq:half}
\end{equation}

\begin{equation}
\begin{cases}
\tilde{x}_m^{(2t)} = x_m^{(2t)} \cos(m\theta_t) - x_m^{(2t+1)} \sin(m\theta_t) \\
\tilde{x}_m^{(2t+1)} = x_m^{(2t+1)} \cos(m\theta_t) + x_m^{(2t)} \sin(m\theta_t)
\end{cases}\label{eq:consec}
\end{equation}
where \( 0 \leq t < d_h/2 \), \(m\) is the position, and the frequency basis \( \theta_t = 10000^{-2t/d_h} \).
As illustrated in Fig. \ref{fig:Rope}, the canonical LLaMA architecture~\cite{HuggingFaceLLaMAGithub} employs the interleaved pairing from eq. \ref{eq:half} . Although this pattern aligns well with the highly parallel and flexible memory systems of GPUs, it is fundamentally mismatched with streaming hardware architectures. In our FPGA-based streaming design, data flows through a function-level pipeline via FIFOs. The non-contiguous access pattern required by the interleaved pairing would force the pipeline to stall while buffering and gathering all elements for a given head.
By contrast, the consecutive pairing from eq. \ref{eq:consec} applied to data within a packed-vector is more hardware-friendly. This approach simplifies data access, reduces FIFO dependencies, and enables a lower initiation interval. To preserve mathematical equivalence with the LLaMA-style implementation while leveraging the benefits of consecutive access, a lossless transformation is necessary for query and key weight matrices. This transformation, applied on a per-head basis, follows the index exchange formula below, which converts the weight organization from the interleaved to the consecutive pattern:
\begin{align}
\begin{cases}
\mathbf{w}_{N_h}^{(2t)} \Leftrightarrow \mathbf{w}_{N_h}^{(t)} \\
\mathbf{w}_{N_h}^{(2t+1)} \Leftrightarrow \mathbf{w}_{N_h}^{(d_h/2 + t)}
\end{cases}
\quad \text{for } 0 \leq t < \frac{d_h}{2}.
\end{align}


% For rope operation, we assume that rope cache is already attained, as computing floating point, exponentiation, sin and cos are time and resource cosuming. Moreover, in the context of LLaMA-alike rope operation, the rope was exerted between the first-half of the hidden dimension and second-half of the hidden dimension. They claim it is more efficient for GPU operation as there is no need to extract the consecutive data as datapairs and map it into a complex domain. However, for the streaming FPGA architecture, consecutive access of the packed vector is more practical and get rid of the memory dependency is fewer initial interval. The Fig. \ref{fig:TLMM_FUSE}(c) displays the two computation implementations and the schedule. Moreover, to ensure the outcome stays the same, the transformation for the weight of the LLaMA-like rope to the original rope is given as follows in the head-wise manner in the following index exchange formula.

\subsection{Weight Buffer Management Unit (WBMU)} \label{SEC:WBMU}
%\vspace{-1mm}
\subsubsection{Analytical Parameter Selection for TLMM based on Optimal URAM Utilization} \label{SEC:PARAMETER_WBMU}
% In the context of LLM accelerator design, the design of external memory access of weight is vital, as they consume most of the time when decoding. Moreover, the on-chip memory is limited in edge FPGAs, and the URAM resource with a capable depth for weight buffering is more limited. At the same time, when the buffer cannot stores the required weight, the performance will be degraded as multiple axi memory access is generated and a partial sum of intermediate results can consume much memory, requiring regenerating a new table when implementing TLMM.
% Therefore, an analytic method for optimal URAM utilization is introduced when designing the on-chip buffer unit of TL-based matrix multiplication.

In the design of LLM accelerators, external memory access for weights is a critical performance bottleneck, often dominating the latency of the decoding phase. This challenge is acute on edge FPGAs, where on-chip memory resources are limited. Among these resources, URAM is ideal for buffering large weight tensors.
Insufficient on-chip buffer capacity for weights leads to significant performance degradation. When the required weights cannot be stored on-chip URAM, the system must repeatedly access off-chip memory via the AXI bus, which will be severely limited by the low DDR bandwidth of the edge FPGA. 
Therefore, this work proposes an analytical model for optimizing URAM utilization by selecting the relevant TLMM parameters. This model provides a systematic approach to designing the on-chip weight buffer and TLMM accelerators, ensuring optimal parameter selection for both modules within the memory constraints of edge devices.
% \subsubsection{Weight Packing Strategy and TLMM Parameter Selection for Maximum URAM Utilization}

As introduced by the \cite{AMD_UG573_ultrascale_memory}, the size of the URAM resource is set as 72 bitwidth $\times$ 4096 depth, which can equal up to 16 BRAM18K blocks. However, due to its fixed bitwidth of 72, the design of storage architecture is vital. We assume there are $N_{\text{URAM}}$ URAM blocks in total in an edge FPGA. The total URAM memory size is $288Kb \times N_{\text{URAM}}$. In our design, each of the units saves a weight index vector $\bf{w}_{idx} \in \bf{W}_{idx}$ and the bit width is $T \times B_{\text{idx}}$. On one hand, our goal is to have a TL table as large as possible within the resource limitations; on the other hand we want to make sure the bitwidth of URAM is fully utilized. Thus, the largest number of TL table $T$ is set according to 
\begin{align}
T = \left\lceil \frac{72 \times c_{\text{URAM}}}{B_{\text{idx}}} \right\rceil
\label{eq:table},
\end{align}
where the factor $c_{\text{URAM}}$ stands for the cascade factor that combine $c_{\text{URAM}}$ URAM together to get a larger bitwidth. The $T$ is also confined by the available LUT resource to ensure place and route to implement the TLMM. The number of \text{LUT} can be utilized for TLMM is $\text{LUT}_{\text{max}}$. According to Sec. \ref{sec:TLMM_METHODS}, the constraint can be given by
\begin{equation}
\begin{split}
   T \times (N_{\text{TB}} \times \text{LUT}_{\text{tree}} + Q &\times N_{\text{TB}} \times 
   \text{LUT}_{\text{entry}} \\&
    + Q \times \text{LUT}_{\text{lp}} )  \leq \text{LUT}_{\text{max}},
\end{split}
\end{equation}
% where $S_{\text{TB}}$ is the size of all $T$ $Q$-duplicated TL tables in bit.
Moreover, to enable the parallel read of $Q$  trait, we will try to partition the target buffer array size of $X_{\text{URAM}} \times Y_{\text{URAM}}$ to ensure the above memory banks are independent.
Since each URAM is a dual-port memory ~\cite{AMD_UG573_ultrascale_memory}, we partition the weight matrix cyclically with a factor of $Q/2$. This strategy ensures that $Q$ consecutive elements are placed in independent blocks, enabling parallel access. 
These logical blocks are then mapped to physical URAMs based on the memory depth. The total number of utilized URAMs $U$ can be given as:
\begin{align}
U = \left\lceil \frac{(\frac {Y_{\text{URAM}}}{  Q} )\times \lceil \frac{X_{\text{URAM}}}{G \times T}\rceil}{4096}  \right\rceil  \times c_{\text{URAM}}  \times \frac{Q}{2}, \quad U \leq N_{\text{URAM}},
\end{align}
where the term wrapped by the ceil gives the depth required in total for a single partitioned array, $c_{\text{URAM}}$ is the cascade factor to build the required bitwidth, and there are $Q/2$ blocks in total. 
% Therefore, according to the constraint rules, the factor $Q$ can be approxiamtely bounded by,
% \begin{align}
% n \leq \frac{E \times depth}{(\frac {S \times Y_{\text{URAM}}}{  Q} )\times \lceil \frac{X_{\text{URAM}}}{G \times T}\rceil} \\
% \frac{T}{n} \approx \left \lceil \frac{bitwidth}{B_{\text{idx}}} \right \rceil \leq \frac{(\frac {S \times Y_{\text{URAM}}}{  Q} )\times \lceil \frac{X_{\text{URAM}}}{G }\rceil}{E \times depth}
% \end{align}
%  Thus, the max group size $G$ can also be approxiamtely obtained according to eq. \ref{eq:table} as follow,
% \begin{align}
% \frac{B_{\text{idx}}}{G}\leq \left\lceil\frac{\frac{bitwidth}{B_{\text{idx}}}\times E \times depth}{(\frac {S \times Y_{\text{URAM}}}{  Q} )\times X_{\text{URAM}}} \right \rceil.
% \end{align}
% where $B_{\text{idx}}=\lceil log_2(3^{G}) \rceil$

% \begin{remark}
% Note that, if $ n \leq 1$, the onchip URAM is not big enough to buffer, the weight loading have to be partitioned into multiple off-chip dram access and the TLMM will also be further tiled to support the model size, leading to in-efficient off-chip memory access.
% \end{remark}
% \begin{figure}
%     \centering
%     \includegraphics[width=\linewidth]{Figures/WBMU_access_management.pdf}
%     \caption{The access pattern of WBMU unit}
%     \label{fig:WBMU_ACCESS}
% \end{figure}
\subsubsection{TLMM Weight Accessing}
\begin{figure}[t]
    \centering
     \includegraphics[width=\linewidth]{Figures/WBMU_SYSTEM_Ye.pdf}
     \vspace{-2mm}
    \caption{The WBMU and AXI data packing strategy. The green arrows denote the connection to the weight access, and the yellow arrows are connected to the weight load module.}
    \label{fig:WBMU_SYSTEM}
\end{figure}
% \todo[inline]{Rectify the color in Fig. 6}

% The 3 main categories of weight sizes in classical LLM: QKVO projection: $d_{\text{model}} \times d_{\text{model}}$ Up projection in FFN $d_{\text{model}} \times d_{\text{ffn}}$, and down projection in FFN $d_{\text{ffn}} \times d_{\text{model}}$. Note that the $d_{\text{model}}$ and $d_{\text{ffn}}$ will be padded to the multiple of $T \times G$ as $d_{\text{model}}' = \lceil \frac{d_{\text{model}}}{T\times G} \rceil \times T \times G$ and $d_{\text{ffn}}' = \lceil \frac{d_{\text{ffn}}}{T\times G} \rceil \times T \times G$
% Now the challenge comes with how to support all TLMM sizes presented in \ref{SEC:TLMM Design}. The first 2 type of matrix multiplication share the common $n = d_{\text{model}}$, so the only change lies in controlling the access of $k$ dimension when doing TLMM. However, Note that the weight size of up and down projection are transposes of each other. More importantly, $d_{\text{ffn}}$ normally does not divide evenly into $d_{\text{model}}$. Therefore, if we want to access the $d_{\text{ffn}}$, we can not access them aligningly, leading to complex none aligned accessing and loading logic, as shown in \todo[inline]{Fig of WBMU}. As a result, we pad the required weight buffer size to make it divisible, so that they can be accessed aligned.
A key challenge is supporting the various TLMM dimensions presented in Sec. \ref{sec:TLMM design}. The first two matmul types share the dimension $n = d_{\text{model}}$, so adapting to them primarily involves controlling memory access along the $k$-dimension. 
However, a complication arises because the weight matrices for the up- and down-projections are transposes of each other, and the address stride is supposed to be different. Furthermore, the dimension $d_{\text{ffn}}$ is not necessarily an even multiple of $d_{\text{model}}$. This dimensional mismatch prevents aligned memory access, which would necessitate complex logic to handle unaligned weight indices access. To resolve this, we pad the weight buffers to make their dimensions evenly divisible, thereby ensuring all memory accesses are aligned.
% Fig. \ref{fig:WBMU_ACCESS} illustrates the weight access pattern for various TLMM sizes in our design, configured with parameters $Q=4$ and $T=4$. To ensure aligned memory access, the last two columns of the weight matrix are padded. 
Moreover, the logical dimensions $d_{\text{ffn}}$ and $d_{\text{model}}$ are padded to be multiples of $T \times G$, resulting in the dimensions $d_{\text{ffn}}'$ and $d_{\text{model}}'$. Therefore, $X_{\text{URAM}} = d'_{\text{FFN}}$ and $Y_{\text{URAM}} = d'_{\text{model}}$.
The WBMU contains an address translation unit that converts a straightforward 2D software weight matrix index $(a, b)$ into a 3D-indexed physical address $(x, y, z)$. This scheme treats the cyclically partitioned URAM array as a single, contiguous address space for simplified indexing explanation. Its array size is configured as [$\left \lceil \frac{d_{\text{ffn}}'}{d_{\text{model}}'} \right \rceil$, $\frac{d_{\text{model}}'}{(T\times G)}$, $d_{\text{model}}$]. With weight index request $1 \leq a,b\leq d_{\text{ffn}}$ from TLMM engine, the address mapping is given by:
\begin{equation}
% \vspace{-2mm}
\begin{cases}
x = \begin{cases}
        0, &\text{q, k, v, o projection}\\
        \frac{b}{d_{\text{model}}} , &\text{Up projection} \\
        \frac{a}{  d_{\text{model}}'}, &\text{Down projection} \\
    \end{cases}\\

y = a \mod  d_{\text{model}}', \quad
z = b \mod d_{\text{model}}, \\
\end{cases}.
\end{equation}

% Fig \ref{fig:WBMU_ACCESS} showcases the weight access pattern of different TLMM sizes in our design when $Q = 4, T=4$ and last 2 weight columns are padded value for aligned access. 
% Thus, the size of the required weight buffer is $X_{\text{URAM}} \times Y_{\text{URAM}}$ is set as 
% $X_{\text{URAM}} =d_{\text{ffn}}'$, 
% $Y_{\text{URAM}} = \left \lceil \frac{d_{\text{ffn}}'}{d_{\text{model}}'} \right \rceil \times d_{\text{model}}'$ and the 
% $d_{\text{ffn}}'$ and $d_{\text{model}}'$ are padded as multiple of $T \times G$.
% Furthermore, the weight accessing unit from WBMU can convert the straightforeward weight matrix index $(a, b)$ into the physical address of the URAM blocks where the URAM blocks are deemed as a whole with cyclic partition style for simple addressing in three dimension $(x, y, z)$. Where $x$ is indexing the $\frac{d_{\text{ffn}}'}{d_{\text{model}}'}$, $y$ is indexing the $d_{\text{ffn}}'$ and $z$ is indexing $d_{\text{model}}'$ Thus the address translating of WBMU weight access is given by,

\subsubsection{Off-chip DDR Weight Transfer}

As illustrated in Fig. \ref{fig:WBMU_SYSTEM}, the WBMU connects to the DDR memory system via three high performance (HP) ports: HP0, HP1, and HP3. The AXI interface for these ports is configured with a 256-bit data width, a burst size of 16, and an outstanding transaction limit of 16, following the guidelines in ~\cite{Benchmark_AXI_ON_ZYNQ}.
During weight loading, the three HP ports transmit data in parallel, as each is connected to an independent DDR Quality of Service (QoS) channel as shown in Fig. \ref{fig:sys_arch}. This parallelism results in a combined bandwidth of 768 bits per AXI access cycle. As shown in Fig. \ref{fig:WBMU_SYSTEM}(b), each 768-bit transfer is structured to carry $\lceil (768 - 16) / (T \times B)\rceil$ ${\bf w}_{\text{idx}}$ and one FP16 RMSNorm weight. In our design, with parameters $T=28$ and $B=5$, a single pipelined burst request can convey up to five weight index vectors and one RMSNorm weight in one transfer out of 16 burst AXI data. 
% The method is supported by the ablation study in \ref{sec:ABLATION_WBMU}.

\subsection{RMS-MAX Unit}
%\vspace{-1mm}
The RMS-MAX hardware unit performs RMS normalization followed by channel-wise maximum finding for subsequent quantization and dequantization. Input data undergoes RMSnorm accumulation with upcasting to FP32 for precision, succeeded by FP16 division and FP16 RMSnorm weight scaling to compute the norm. A stream FIFO channel interfaces the RMSnorm output with the channel-wise max finding module: vector-wise maximums are first extracted per channel segment, and upon completion of the final channel flow, the global channel maximum is identified and stored in a dedicated buffer for quant/dequant operations. By decoupling the max-finding and quantization logic, memory dependencies are eliminated in computing the final quantized output. Furthermore, max-finding can proceed concurrently with vector-wise output FIFO due to their strictly sequential nature.
% From the WBMU system in Fig. \ref{}(a), the weight URAM buffer is connected to 3 HP ports- HP 0, HP1 HP3. The AXI burst size of 16, AXI outstanding size of 16, and an AXI bitwidth of 256, which follows the instruction in ~\cite{Benchmark_AXI_ON_ZYNQ}. When the weight is loading, the three HP ports will transmit the data in parallel because they are connected to an independent DDR QoS channel. In total, there are $768$ bit transferred at one parallel AXI access of the three ports. As shown in Fig. \ref{}(b), at one access, $\lceil (768 - 16) / (T \times B)\rceil$ weight index vectors ${\bf w}_{\text{idx}}$ and one FP16 rms-weight will be transfered. In our design, the $T$ is set at 28 an $B$ is 5. Thus, up to 5 weight index vectors and one rmsnorm can be transferred, with a pipelined request of burst transfer. 
% \todo[inline]{Describe the way of 3port transporting multiple data}


% \subsubsection{TLMM Weight Accessing}
% After loading, when running TLMM, the weight will be accessed in 2 different manners in different matrix multiplication manner. When $n = d_{\text{model}}$, the access pattern is sequential. Moreover the RMS weights sized $1 \times d_{\text{model}}$ and $1 \times d_{\text{ffn}}$ are also sent to the on-chip RMS weight buffer along with the indices weight. The utilization of the AXI bitwidth will be $B_{\text{idx}} \times T + 16$ as the RMS weight is represented in FP16. 
% As for the down projection $n = d_{\text{ffn}}$ in the previously mentioned part,the access will be combine the access of $\left\lceil \frac{d_{{ffn}}}{d_{\text{model}}} \right\rceil$ none consecutive memory region to get the corresponding size. The details of accessing address generation is given in the \todo[inline]{pseudocod for address access mapping}

% Assume the weight buffer is represented by a arrary WeightThe detailed address mapping is given as follows,
% \begin{align}

% \end{align}






% \section{RMS-MAX Unit}
% \subsection{detailed design of RMS-MAX Engine}
% \subsection{Analysis for Design Choice of Operator Fusion in RMS-MAX and TLMM-FUSE}

% \todo[inline]{RMS MAX PENDING}
% \section{Specialized Reverse Reorder for Prefill Attention}
% \section{Attention in Prefill And Decode}
% %\vspace{-1mm}
% \usepackage{subcaption}
% \usepackage{graphicx}
% \usepackage{captionsetup} % if you already load 'caption', you can skip this line
\captionsetup{font=small}
\captionsetup[subfigure]{font=footnotesize,labelformat=parens,labelsep=space}

% ---------------------------------------------------------
\begin{figure*}[t]
\centering
\setlength{\abovecaptionskip}{4pt}
\setlength{\belowcaptionskip}{-4pt}
\setlength{\floatsep}{6pt}
\setlength{\textfloatsep}{6pt}
\begin{subfigure}{0.35\textwidth}
  \centering
  \includegraphics[width=\linewidth]{Figures/reverse_overview.pdf}
  \caption{Schedule on the Attention Score Map ($N_{\text{pe}}{=}4$)}
  \label{fig:attention_map}
\end{subfigure}\hfill
\begin{subfigure}{0.29\textwidth}
  \centering
  \includegraphics[width=\linewidth]{Figures/normal_schedule_v2.pdf}
  \caption{Naive Attention Scheduling ($N_{\text{pe}}{=}4$)}
  \label{fig:naive_schedule}
\end{subfigure}\hfill
\begin{subfigure}{0.35\textwidth}
  \centering
  \includegraphics[width=\linewidth]{Figures/reverse_schedule_v2.pdf}
  \caption{Reverse Attention Scheduling ($N_{\text{pe}}{=}4$)}
  \label{fig:reverse_schedule}
\end{subfigure}
 \caption{Attention Schedules Comparison}
\label{fig:attention_triptych}
\end{figure*}


% \begin{figure}[t] 
% \centering
% \includegraphics[width=\linewidth]{Figures/reverse_overview_Ye.pdf}
% \vspace{-3mm}
% \caption{The Schedule on the Attention Score Map ($N_{\text{pe}} = 4$, beige stands for attention mask)}
% \label{fig:attention_map} 
% \end{figure}



% \begin{figure}[t] 
% \centering
% \includegraphics[width=\linewidth]{Figures/naive_schedule_Ye.pdf}
% \vspace{-3mm}
% \caption{Naive Attention Scheduling ($N_{\text{pe}} = 4$)}
% \label{fig:naive_schedule} 
% \end{figure}

% \begin{figure}[t] 
% \centering
% \includegraphics[width=\linewidth]{Figures/reverse_Ye.pdf}
% \vspace{-3mm}
% \caption{Reverse Attention Scheduling ($N_{\text{pe}} = 4$)}
% \label{fig:reverse_schedule}
% \end{figure}


\subsection{Reversed Prefill Attention Unit}
%\vspace{-1mm}
\subsubsection{Prefill Attention Challenge on edge FPGAs}
% Prefill is one of the most tricky part for the edge FPGA, this part of LLM requires abundant resource and bandwidth for multi-token computation, especially the attention computation that requires softmax and the matrix to matrix multi-head operations with the complexity of $n^2$.  With limited memory bandwidth and finite computation unit, the computation order of prefill attention should be specially scheduled to meet the requirements. Otherwise, it will be restricted by the edge FPGA bandwidth. In the field of edge FPGA vision transformer, ~\cite{edgemoe} has proposed the state-of-the-art dense attention scheduling as depicted in Fig. \ref{fig:attention_map} that considers the reuse of Q value within different token. However, different from the vision transformer, the causal attention mask is applied for LLMs. Thus, the dense scheduling wastes computation cores on the zero masks.  
% Moreover, operation fusion of $Q\otimes k$, $softmax$ and $score \otimes V$ can be applied to reduce the extra access of DDR. The state-of-art kernel fusion implementation on resource-abundant GPU is the flash attention. Nevertheless, the GPU adapted computation is not suitable for FPGA, because it has many more computation cores than FPGA and the on-chip SRAM is much larger than on-chip BRAM/URAM of FPGA. Therefore, we proposed the \textit{reverse attention} utilized the fused attention and reverse reorder scheduling, specialized for edge FPGA.
The prefill attention of LLM requires significant resources and bandwidth for multi-token computation at the sequence length of $N$, especially for attention computation, which involves softmax and matrix-to-matrix multi-head operations with a complexity of \( N^2 \). Given the limited memory bandwidth and finite computational units of DSP for FP operation on edge FPGAs, the computation order of prefill attention should be carefully scheduled to meet these requirements. Otherwise, it may be constrained by the bandwidth limitations of the edge FPGA, as shown in the naive attention scheduling in Fig. \ref{fig:naive_schedule} \cite{edgemoe}. When we have $N_{\text{pe}}$ processing elements (PEs), the bandwidth requirement will be $N_{\text{pe}}$ as well. The more PE we have, the larger the amount of data will be required from DDR, bounded by the total bandwidth. 
%as the size of the intermediate softmax score matrix is $L \times L \times head$. 
%as shown in the naive attention scheduling in Fig. \ref{fig:naive_schedule}. 

% In the context of edge FPGA vision transformers, ~\cite{edgemoe} proposed a state-of-the-art dense attention scheduling strategy, as shown in Fig. \ref{fig:dense_schedule} and Fig. \ref{fig:attention_map}, which takes into account the reuse of the \( \bf Q \) vector across different tokens. However, unlike vision transformers, LLMs use a causal attention mask. As a result, the dense scheduling approach wastes computational resources on zero masks.

Furthermore, $N \times N \times N_{\text{head}}$ elements are stored back to DDR and loaded on chip as the softmax matrix $\bf S$.
The fusion of operations such as \( {\bf Q} {\bf K}^{\text{T}} \), softmax, and \( \bf S  {\bf V} \) can reduce the additional accesses to DDR. The state-of-the-art kernel fusion implementation for resource-abundant GPUs is Flash Attention \cite{dao2022flashattentionfastmemoryefficientexact}. However, GPU-optimized computation is not suitable for FPGAs, as GPUs have many more computational cores and much larger on-chip SRAM compared to the on-chip BRAM/URAM available on FPGAs. To address these challenges, we propose the reverse attention method, which utilizes fused attention and reverse reorder scheduling, specifically tailored for edge FPGAs.

% \begin{table}[ht]
% \caption{Comparison of different attention approaches.}
% \centering
% \resizebox{\linewidth}{!}{
% \begin{tabular}{|c|c|c|c|c|}
% \hline
% \textbf{Approach} & \textbf{KV Data Block Load} & \textbf{Iteration Count} & \textbf{Bandwidth} \\
% \hline
% Our Design  & $\frac{n^2}{2p} + \frac{n}{2})  $ & $\frac{n^2}{2p} + \frac{n}{2}$ & $\sim 1$ \\
% \hline
% Naive & $n^2 + n$ & $\frac{n^2}{p}$ & $\sim p$ \\
% % \hline
% % dense scheduling ~\cite{edgemoe}& $\frac{n^2}{p} + n + p - 1$ & $\frac{n^2}{p} + p - 1$ & $\sim 1$ \\

% \hline
% \end{tabular}
% }
% \label{tab:compare}
% \end{table}

% \begin{figure}[t] 
% \centering
% \includegraphics[width=0.85\linewidth]{Figures/dense_schedule.pdf}
% \caption{Dense attention scheduling ($p = 4$).}
% \label{fig:dense_schedule} 
% \end{figure}


\begin{figure}[t] 
\centering
\includegraphics[width=\linewidth]{Figures/reverse_attention_PE_arch_Ye.pdf}
\caption{Design of Fused Attention at Iteration Step 4 ($N_{\text{PE}} = 4$)}
\label{fig:reverse_pe} 
\vspace{1em}
\end{figure}

 \begin{figure}[ht]
    \centering

    \includegraphics[width=\linewidth]{Figures/decoding_attention.pdf}

    %\caption{The Structured Data Access for SpMV}
    \caption{Detailed Design of Decoding Attention}
    %%\Description{}
    \label{fig:decoding_atten}
\end{figure}
\subsubsection{Reversed and Fused Attention}
% The pseudo-code of \textit{reverse attention} is shown in Algorithm \todo{algorithm} and the scheduling is shown in Fig. \ref{fig:reverse_schedule}. Assume that the current length of prefill tokens is $n$. $1< i \leq n$ and $1 < j \leq n$ indicate the current token index for $\bf q$ and $\bf k$, $\bf v$. There are $N_{h}$ heads in total.

% The kernel fusion computation can be deemed a special case for Flash Attention V2 ~\cite{dao2022flashattentionfastmemoryefficientexact} when block size equals 1. The head-wise formula of the case of 2 consecutive block can be written as,  

The reverse attention scheduling is depicted in Fig. \ref{fig:reverse_schedule}. Assume that the current sequence length of the prefill tokens is \( N \), with \( 1 < i \leq N \) and \( 1 < j \leq N \) representing the current token indices for \( {\bf q} \) and \( {\bf k}, {\bf v} \), respectively. There are a total of \( N_{{\text h}} \) heads. The operator of $\bf{qkv}$ denotes the fused operation.
The kernel fusion computation can be considered a special case of Flash Attention V2 ~\cite{dao2022flashattentionfastmemoryefficientexact} when the block size is equal to 1. The head-wise formula for the case with two consecutive blocks can be written as follows:
\begin{equation}
\vspace{-2mm}
\left\{
\begin{small}
\begin{aligned}
  s^{(i)} &= \frac{ {\bf q}_{i}{\bf k}^{\text{T}}_{i}}{\sqrt{d_{\text{h}}}}, m^{(i)} = \max\left(m^{(i-1)}, s^{(i)}\right)\\
\ell^{(i)} &= e^{m^{(i-1)} - m^{(i)}} \ell^{(i-1)} + e^{s^{(i)} - m^{(i)}} \\
\bf{o}^{(i)} &= e^{m^{(i-1)} - m^{(i)}}{ \bf o}^{(i-1)} + e^{s^{(i)} - m^{(i)}} {\bf v}_{i}
\end{aligned}
\end{small}
\right.
% \vspace{-0.8mm}
\end{equation}
where $m^{(i)}$ denotes the running maximum value, $s^{(i)}$ represents the dot product score, $\ell^{(i)}$ is the denominator factor, and ${\bf o}^{(i)}$ is the numerator vector. After all inputs have been processed (up to step N), the final attention output is computed as the division of the final numerator by the final denominator ${\bf o}^{(i)}/\ell^{(i)}$. 
Compared to a naive implementation, this online algorithm avoids generating the large intermediate attention score matrix $\bf S$, saving significant off-chip DDR access. It also breaks the data dependency on finding a global maximum value, which enables streaming computation and operation fusion.

% The detailed hardware detailed design is depicted in Fig. \ref{fig:reverse_pe}.  

%There are $N_{h}$ of these intermediate vectors in total.

% \begin{equation}
% \left\{
% \begin{aligned}
% m^{(1)} &= s^{(1)} \\
% \ell^{(1)} &= e^{s^{(1)}-m^{(1)}} \\
%{ \bf o}^{(1)} &= e^{s^{(1)}-m^{(1)}} {\bf v}^{(1)} \\
% m^{(2)} &= \max \left(m^{(1)}, s^{(2)}\right) = m \\
% \ell^{(2)} &= e^{m^{(1)}-m^{(2)}} \ell^{(1)} + e^{s^{(2)}-m^{(2)}} \\
% &= e^{s^{(1)}-m} + e^{s^{(2)}-m} = \ell \\
% \bf{p}^{(2)} &= e^{s^{(2)}-m^{(2)}} / \ell^{(2)} \\
%{ \bf o}^{(2)} &={ \bf o}^{(1)} / e^{m^{(1)}-m^{(2)}} + e^{s^{(2)}-m^{(2)}} {\bf v}^{(2)} \\
% &= e^{s^{(1)}-m} {\bf v}^{(1)} + e^{s^{(2)}-m} {\bf v}^{(2)} \\
%{ \bf o}^{(2)} &={ \bf o}^{(2)} / \ell^{(2)} ={ \bf o}
% \end{aligned}
% \right.
% \tag{3}
% \end{equation}
% where \( m \) denotes the maximum value, \( s = {\bf q}_i \otimes {\bf k}_i \) represents the MAC result of the vector dot product, \( \ell \) is the denominator factor, and \({ \bf o} \) is the numerator vector. The upper index indicates the current step in the kernel fusion computation.

Regarding detailed schedule, instead of starting from the first token \( {\bf q}_1 \), our schedule begins from \( {\bf q}_{N-1} \). Specifically, the level of parallelism is set to \( N_{\text{pe}} \). The factor \( N_{\text{pe}} \) also implies that the PE buffer can store \( N_{\text{pe}} \) tokens of \( {\bf q}_i \). In each iteration, one \( {\bf q}_i \) token is loaded onto the buffer of PEs with unicast (with the first batch loading \( {\bf q}_N \) to \( {\bf q}_{N-3} \)). Simultaneously, the corresponding \( {\bf k}_j \) and \( {\bf v}_j \) tokens are loaded to PEs in multicast. After all \( N \) \( {\bf k}_j \) and \( {\bf v}_j \) tokens have been loaded and the fused-kernel computation is completed, the next iteration will evict \( N_{\text{pe}} \) \( {\bf k}_j \) and \( {\bf v}_j \) tokens, starting from \( {\bf k}_{N-3} \) and \( {\bf v}_{N-3} \), to avoid redundant computations arising from the causal attention mask. Once all computation with regard to the loaded $\bf q$s is completed, the  ${\bf o}^{(i)}/\ell^{(i)}, i = N \text{ to } N-3$ for all heads are pushed to the as consecutive heads consist the hidden dimension. The microarchitecture state at step 4 is depicted in Fig.~\ref{fig:reverse_pe}. Each PE is divided into a head-wise multiplication and accumulation (MAC) unit, a rescale and accumulation unit for the output vector and denominator factor, and a final division and output unit.

The iteration continues until all \( 1 < i \leq N \), \( 1 < j \leq N \) are traversed. In this approach, the only required input buffers are for \( N_{\text{pe}} \) \( {\bf q}_i \) tokens, one \( {\bf k}_j \), and one \( {\bf v}_j \). Additionally, the intermediate buffers of the $N_{h}$ heads include: \( N_{h} \times N_{\text{pe}} \) multi-head MAC intermediate results \( s \), \( N_{h} \times N_{\text{pe}} \) multi-head previous max values \( m \), \( N_{h} \times N_{\text{pe}} \) intermediate denominators \( \ell \), as well as the output vector $\bf o$ buffer sized $d_{\text{model}} \times N_{\text{pe}}$. The trade-off also includes minor computational overhead: an additional exponentiation, $e^{m^{(i-1)} - m^{(i)}}$, is required to rescale previous results when the running maximum is updated. Additionally, the state of the running numerator $\bf{o}$ and denominator $\ell$ are maintained with two extra floating point rescales. 

To support the reverse reorder in the prefill stage, the embedding vectors are indexed by the reversed input prompt reordered token id, and the rope cache is also reversed. In this way, the RPA can be seamlessly utilized without extra reversed address DDR access logic that is not supported by the AXI address incremental burst access protocol \cite{amd_ug1037}. The advantage of the reversed attention is further proved in Sec. \ref{SEC:ABLATION_ATTENTION}.
% After introducing the reverse attention process, a comparison between naive attention and our proposed reverse attention is provided in Table \ref{tab:compare}. The comparison indicates that the iteration count for reverse attention is the lowest, and the required bandwidth remains constant. Furthermore, the redundant outcome rate of the computation core can be directly assessed from the attention map in Fig. \ref{fig:attention_map}. The naive scheduling approaches do not account for the causal mask, leading to many redundant computed values. 
%It is important to note that due to kernel fusion, latency tests cannot be directly compared, and results will be presented in \todo{ablation study}.
% where $m$ indicates the max value, $s = q_{i} \otimes k_{i}$ is the MAC result of vector dot product, $\ell$ is the denominator factor, $\bf o$ is the numerator vector. The upper index indicates the current step of kernel fusion computation.
% As for scheduling, instead of starting from the first token $\bf q_{1}$, our schedule starts from the ${\bf q}_{n-1}$. In detail, the level of parallelism is set as $p$ (the figure illustrate the situation of $p = 4$). The $p$ factor also indicates that on-chip BRAM can stores $p$ $\bf q_{i}$ tokens. For each iteration, one $\bf q_{i}$ token is loaded to the onchip (the first batch will be ${\bf q}_{n}$ to ${\bf q}_{n-3}$). Meanwhile, for every iteration, the corresponding ${\bf k}_{j}$ and ${\bf v}_{j}$ are loaded for computation. After all $n$ ${\bf k}_{j}$ ${\bf v}_{j}$ tokens are loaded and finished fused-kernel computation, next iteration will evict p ${\bf k}_{j}$ ${\bf v}_{j}$ and start from ${\bf k}_{n - 3}$ ${\bf v}_{n-3}$ to avoid redundant computation that stems from the casual attention mask. The iteration carries on until all $1 < i \leq n$, $1 < j \leq n$ are traversed. In this way, the only required input buffers are for $p$ ${\bf q}_{i}$, one ${\bf k}_{j}$ and one ${\bf v}_{j}$. Moreover, the intermediate buffers are: $N_{h} \times p$ multi-head MAC intermediate results $s$, $N_{h} \times p$ multi-head previous max value $m$, and $N_{h} \times p$ intermediate denominator $\ell$.

% After introducing the process of reverse attention, a comparison between naive attention, dense attention in Edge Moe ~\cite{edgemoe} and our proposed reverse attention can be given in \ref{tab:compare}. The comparison indicates that the iteration count for reverse attention is the lowest, and the required bandwidth is constant. Moreover, the redundant outcome rate of the computation core can be directly compared from the attention map in Fig. \ref{fig:attention_map}. The naive and dense scheduling didn't consider the causal mask and many computed values are redundant. Notice that due to the kernel fusion, latency test can not be directly compared and will be displayed in \todo{ablation study}. 
% Table \todo{Bandwidth/IDLE RATE analysis} displays the general case

%\subsection{Unified Decoding Attention and LM-Head}



\subsection{Decode Attention Unit}
%\vspace{-1mm}
The computational characteristics of attention differ between prefill and decoding phases. The computation in the decoding phase attention involves primarily matrix-vector and vector-vector operations. The computational load per step $O(Nd_{\text{model}})$ is significantly lower than the total matmul in prefill attention, $O(N^2d_{\text{model}})$. However, this phase requires fetching the large ${\bf K}_{\text{cache}}$ and ${\bf V}_{\text{cache}}$ matrices from memory (e.g., off-chip DDR) in every layer. Consequently, the decoding phase is often memory-bandwidth bound, especially as the sequence length $N$ grows. 
This heterogeneity of these workloads necessitates a system-level balance between performance and resource utilization to achieve optimal end-to-end efficiency.
For decoding attention, the marginal benefits of adding extra compute resources diminish compared with prefill attention. To address this, we decouple the fusion between $\bf{QK^{\text T}}$ multiplication (fetch ${\bf K}_{cache}$) and the weighted aggregation with ${\bf V}$ (fetch ${\bf V}_{cache}$), since these two operations are dominated by memory access. 
 Fig. \ref{fig:decoding_atten} illustrates the hardware architecture of our decoding attention design. The query input \textbf{Q} (single token) is pre-loaded on-chip, while the 
${\bf K}_{\text{cache}}$ and ${\bf V}_{\text{cache}}$ is loaded from off-chip DDR via AXI interface. We utilize a stream-like dataflow within the module to conceal the computation latency. We use an online softmax algorithm that partially fuses the softmax with $\bf{QK^{\text T}}$ vector-matrix multiplication. The intermediate attention score is stored in the on-chip stream buffer, so no off-chip memory access is required in softmax computation.


% The softmax computation in the decoding attention does not require off-chip memory access (only one token processed). The softmax score matrix $\bf S$. We use an online softmax algorithm that partially fuses the softmax with QK multiplication. It hides non-memory latency with negligible resource overhead.

% In Sec. , we introduce a prefill attention unit that fully fuses the three major steps of attention computation: QK multiplication, softmax normalization, and the weighted aggregation with V. Since each step in prefill attention involves intensive computation, this fused design significantly reduces latency. However, this fully pipelined implementation requires calculating additional intermediate results, which introduces extra hardware resources.







% In general, the bit-wise operation for mixed-precision matrix multiplication can be expressed as follows:
% \begin{align}
%     {\bf A} \otimes {\bf W} = {\bf A} \otimes \left(\sum_{i = 0}^{n-1} 2^i {\bf W}_{i}\right) = \left(\sum_{i = 0}^{n-1} 2^i {\bf A} \otimes {\bf W}_{i}\right)
% \end{align}

% Considering the ternary matrix multiplication with \( W \in \{-1, 0, 1\} \), the above equation can be simplified to:

% \begin{align}
%     {\bf A} \otimes {\bf W} = {\bf A} \otimes {\bf W}_{0}, \quad {{\bf W}_{\text{ternary}}} \in \{-1, 0, 1\}
% \end{align}
% In this case, simple summation and subtraction can replace the multiplication process. However, the method of selecting -1, 0, and 1 to determine the summation and subtraction may not be optimal, as there are only a limited number of combinations of -1, 0, and 1, leading to repetitive computations for the corresponding \( \bf A \) entries. Furthermore, when increasing computation parallelism by duplicating the selection unit of the adding and subtracting path, the resource consumption of the selection may exceed that of the TL tables themselves. This is because the multiple reading ports of the on-chip distributed RAM unit can support multiple accesses to the TL tables, requiring only additional buffers for addressing \todo[inline]{and buffering}. The supportive ablation study will be presented in the next subsection.

% \begin{algorithm}
% \tiny

% \caption{TL-based Ternary matrix multiplication}
% \label{alg:tmac_stream}
% \SetKwInOut{Input}{Input}
% \SetKwInOut{Output}{Output}

% \Input{
%     $\bf \bf A$: Input activation stream (shape $[m][n]$)\;
%     ${\bf W}_{idx}$ = Offline\_preprocess(${\bf W}$): Offline-preprocessed weight indices  (shape $[n/(T*G)][k]$)\;
% }
% \Output{
%     $\bf O$: Output activation stream (shape $[m][k]$)\;
% }
% \BlankLine
% Initialize:\;
% \Indp
%     $\mathit{TL\_TABLE}[n][3^{G}] \gets 0$ \tcp*[r]{Table for all signed combinations}
%     $\mathit{A\_BLOCK}[T \times G ] \gets 0$ \tcp*[r]{Activation buffer}
%     $\mathit{O\_BLOCK}[k] \gets 0$ \tcp*[r]{Output vector accumulator}
% \Indm

% \For{$i \gets 0$ \KwTo $m - 1$}{
%     \For{$j \gets 0$ \KwTo $n - 1$ \textbf{step} $T \times G$}{
%         \tcp{Load activation block}
%         \For{$p \gets 0$ \KwTo $T \times G - 1$}{
%             $\mathit{ A\_BLOCK}[p] \gets \text{A.read()}$\;
%         }
        
%         \tcp{Set up values of TL\_TABLE}
%         \For{$t \gets 0$ \KwTo $T - 1$}{
%             $\mathit{val}_{1\dots G} \gets \mathit{A\_BLOCK}[t \times G : (t + 1)\times G  -1]$\;
%          \text{TL\_TABLE\_set\_up}$(\mathit{val}_{1\dots G})$\;
%         }
        
%         \tcp{Process hidden dimension}
%         \For{$m \gets 0$ \KwTo k \textbf{step} $Q$}{
%             \For{$n \gets 0$ \KwTo $Q - 1$}{
%                 $\mathit{idx\_vec} \gets B_{\text{idx}}\left[\left\lfloor\frac{j}{T \times G}\right\rfloor\right][m + n]$\;
%                 \For{$t \gets 0$ \KwTo $n - 1$}{
%                     $\mathit{TL\_TABLE\_idx} \gets \mathit{idx\_vec}[t]$\;
%                     $\mathit{O\_BLOCK}[m + n] \gets \mathit{O\_BLOCK}[m + n] + \mathit{TL\_TABLE}[t][\mathit{TL\_TABLE\_idx}]$\;
%                 }
%             }
%         }
%     }
    
%     \tcp{Write output}
%     \For{$p \gets 0$ \KwTo $k - 1$}{
%         $\text{{\bf O}.write}(\mathit{C\_BLOCK}[p])$\;
%         $\mathit{O\_BLOCK}[p] \gets 0$\;
%     }
% }

% \BlankLine
% \SetKwProg{Fn}{Function}{:}{}
% \Fn{\text{Offline\_preprocess}($\bf W$)}{
%     \Return Encode every $G$ value as an index in the matrix and pack every $T$ values as a index\_vector $idx\_vec$\;
% }
% \BlankLine
% \SetKwProg{Fn}{Function}{:}{}
% \Fn{\text{TL\_TABLE\_set\_up}($\mathit{val}_{1\dots G}$)}{
%     \Return return all $3^G$ add and subtract combination\;
% }
% \end{algorithm}




% \begin{figure}[t] 
% \centering
% \includegraphics[width=\linewidth]{Figures/lutmatrix multiplication.pdf}
% \caption{Dataflow and architecture of TL-based ternary matrix multiplication ($G = 4$)}
% \label{fig:lutmatrix multiplication} 
% \end{figure}

% % As introduced in \todo{algorithm} and Fig. \ref{fig:lutmatrix multiplication}, the look up table-based matrix multiplication can be divided into 2 stages: (1) the weight is preprocessed into groups (2)  the online ternary matrix multiplication computation. Before computation. Assume that we pack up every 3 ternary values into one index for look-up table addressing so that we have $3\times 3 \times 3 = 27$ combinations and $\log_{2}27 \approx 5$bit for pack-up index representation. $A \in 8B^{n \times m}$ and $W \in ternary^{m \times k}$. The preprocessing of weight is encoding every $G = 3$ ternary values into the 5-bit pack-up index. Then, in the online precomputation stage, we are doing vector-wise tiled matrix multiplication witch is sized $A \in ternary^{1 \times n}$ and $B_{\text{idx}} \in 8B^{n \times k}$. The first step is establishing the look-up table by the precompute unit, which consists of 27 sets of adder/subtracts for all possible combinations of the dynamic activation. After that, the previous packed-up indexes are used to address and select the corresponding value in the look-up table. Then the values are accumulated accordingly to generate a single vector $O \in 8B^{1 \times k}$ in the output matrix. 


% As described in \textbf{Algorithm} \ref{alg:tmac_stream} and Fig. \ref{fig:lutmatrix multiplication}, the TL-based matrix multiplication can be divided into two stages: (1) preprocessing the weights into groups sized $G$, and (2) performing the online ternary matrix multiplication computation. 

% In the preprocessing stage, assume that every $G = 3$ ternary values are packed into a single index for TL table addressing, resulting in \(3^{G} = 3 \times 3 \times 3 = 27\) combinations. The index representation for this packing requires \( \log_2 27 \approx 5 \) bits. Let \( {\bf A} \in \mathbb{N}^{m \times n} \) and \({\bf W} \in \text{ternary}^{n \times k} \). The preprocessing of the weights involves encoding every group of \( G = 3 \) ternary values into a 5-bit packed index.

% In the online stage, we perform vector-wise tiled matrix multiplication, where \( {\bf A} \) and \( \bf B_{\text{idx}} \). The first step is to establish the TL table using the precompute unit, which consists of $3^{G} = 27$ sets of adders and subtractors to cover all possible combinations of dynamic activations. The packed-up indexes are then used to address and select the corresponding values from the TL table. Finally, the selected values are accumulated to generate a single vector \( {\bf O} \in \mathbb{N}^{1 \times k} \) in the output matrix.



% \subsubsection{Multi-TL table dataflow design}
% % In our design, we carefully design the dataflow of the look up table matrix multiplication as shown in Fig. \ref{fig:lutmatrix multiplication}. Assume, we have $T$ tables in total. Firstly, to better vectorize the Loop-up table matrix multiplication. The continued $T $ indexes can be grouped in an index vector, as they can simultaneously access different Look-up tables (i). The weight index vector matrix can be rewritten as $W_{idx}{ \in \{5B, T\}^{\frac{n}{T} \times k}}$. 

% % As for the scheduling, the innermost loop is composed by first the vector operations of establishing $T$ look-up tables simultaneously according to the first $T \times G$ entries of the $A$ matrices. Then utilizing the NR traits of URAM to utilize $Q$ $W$ index vectors in parallel for look-up table addressing that returns $Q* T$ outcome. Then, accumulate the corresponding look-up table return value to the output buffer sized $k$. The $k$ index vectors on the row of $W$ are traversed with the step of $Q$. The look-up table addressing and accumulation process can be fully pipelined with an interval of 1 cycle because there is no inter-iteration dependency. As the first $T \times G$ entries of the $A$ matrices finish, we traverse the $m$ values of the row of the $A$ with the step of $T \times G$ in the intermediate loop. Lastly, the outmost loop traverse the different row vectors in $A$, i.e. the tokens in the prefill stage of LLM. 


% In our design, we carefully optimize the dataflow for TL-based matrix multiplication, as illustrated in Fig. \ref{fig:lutmatrix multiplication}. Assume we have a total of \( T \) tables. To better vectorize the TL-based matrix multiplication, the consecutive \( T \) indices can be grouped into an index vector, enabling simultaneous access to different look-up tables. The weight index vector matrix can be rewritten as \( W_{\text{idx}} \in \{5B, T\}^{\frac{n}{T\times G} \times k} \). \( W_{\text{idx}}\) are loaded onto the on-chip URAM.

% Regarding the scheduling, the innermost loop first performs vector operations to establish the \( T \) look-up tables simultaneously, based on the first \( T \times G \) entries of the \( A \) matrix. Then, leveraging the multiple reading traits of the URAM ~\cite{URAM_Bind}, \( Q \) \( W \) index vectors are processed in parallel for TL table addressing, returning \( Q \times T \) outcomes. The corresponding TL table return values are then accumulated into an output buffer of size \( k \). The \( k \) index vectors on each row of \( W \) are traversed in steps of \( Q \). The TL table addressing and accumulation process can be fully pipelined with an interval of one cycle, as there are no inter-iteration dependencies. After the first \( T \times G \) entries of the \( A \) matrix are processed, the \( m \) values of the row of \( A \) are traversed in steps of \( T \times G \) in the intermediate loop. Finally, the outermost loop traverses the different row vectors in \( A \), corresponding to the tokens in the prefill stage of the LLM.

% As for comparison, the LUT consumption of different matrix multiplication unit design methods is also presented in Table \ref{tab:lut_comparison}. The configuration for the TL-based matrix multiplication is set as \( G = 3 \), \( T = 32 \), and \( Q = 16 \). All levels of parallelism for the modules are set to be the same. Our design consumes 52094 LUTs, while the naive implementation, which selects whether to add or subtract, requires 7905 more LUTs. Another approach involves storing half of the possible combinations (13 out of 27) instead of all combinations and using the index to determine whether the value should be negative. This approach results in a smaller distributed RAM size, aiming to save LUTs. However, after synthesizing, it consumed 9209 more LUTs.

% \begin{table}[N_{h}]
% \centering
% \caption{LUT Consumption Comparison of Different matrix multiplication Unit Design Methods}
% \resizebox{\linewidth}{!}{%
% \begin{tabular}{|c|c|c|}
% \hline
% \textbf{Approach} & \textbf{LUT Consumption} & \textbf{Difference} \\

% \hline
% TL-based(Our Design) & 52,094 & -- \\
% \hline
% Naive Implementation & 59,999 & +7,905 \\
% \hline
% Partial Storage  & 61,303 & +9,209 \\
% \hline

% \end{tabular}%
% }
% \label{tab:lut_comparison}
% \end{table}



% %\todo{remark on the dataflow scheduling}

% \begin{figure}[t] 
% \centering
% \includegraphics[width=\linewidth]{Figures/system_arch.pdf}
% \caption{System architecture of TeLLMe.}
% \label{fig:sys_arch} 
% \end{figure}



% \subsection{Specialized Reverse Reorder for Prefill Attention}



% \begin{figure*}[t] 
% \centering
% \includegraphics[width=0.85\linewidth]{Figures/reverse.pdf}
% \caption{The visualization of scheduling on the attention map (number of computation core $p = 4$, beige stands for attention mask).}
% \label{fig:attention_map} 
% \end{figure*}



% \begin{figure}[t] 
% \centering
% \includegraphics[width=\linewidth]{Figures/naive_schedule.pdf}
% \caption{Naive attention scheduling ($p = 4$).}
% \label{fig:naive_schedule} 
% \end{figure}




% \subsubsection{Prefill Challenge on edge FPGA}
% % Prefill is one of the most tricky part for the edge FPGA, this part of LLM requires abundant resource and bandwidth for multi-token computation, especially the attention computation that requires softmax and the matrix to matrix multi-head operations with the complexity of $n^2$.  With limited memory bandwidth and finite computation unit, the computation order of prefill attention should be specially scheduled to meet the requirements. Otherwise, it will be restricted by the edge FPGA bandwidth. In the field of edge FPGA vision transformer, ~\cite{edgemoe} has proposed the state-of-the-art dense attention scheduling as depicted in Fig. \ref{fig:attention_map} that considers the reuse of Q value within different token. However, different from the vision transformer, the causal attention mask is applied for LLMs. Thus, the dense scheduling wastes computation cores on the zero masks.  
% % Moreover, operation fusion of $Q\otimes k$, $softmax$ and $score \otimes V$ can be applied to reduce the extra access of DDR. The state-of-art kernel fusion implementation on resource-abundant GPU is the flash attention. Nevertheless, the GPU adapted computation is not suitable for FPGA, because it has many more computation cores than FPGA and the on-chip SRAM is much larger than on-chip BRAM/URAM of FPGA. Therefore, we proposed the \textit{reverse attention} utilized the fused attention and reverse reorder scheduling, specialized for edge FPGA.
% Prefill is one of the most challenging components for edge FPGAs. This part of LLM requires significant resources and bandwidth for multi-token computation, especially for attention computation, which involves softmax and matrix-to-matrix multi-head operations with a complexity of \( n^2 \). Given the limited memory bandwidth and finite computational units, the computation order of prefill attention must be carefully scheduled to meet these requirements. Otherwise, it may be constrained by the bandwidth limitations of the edge FPGA, as shown in the naive attention scheduling in Fig. \ref{fig:naive_schedule}. 

% In the context of edge FPGA vision transformers, ~\cite{edgemoe} proposed a state-of-the-art dense attention scheduling strategy, as shown in Fig. \ref{fig:dense_schedule} and Fig. \ref{fig:attention_map}, which takes into account the reuse of the \( \bf Q \) values across different tokens. However, unlike vision transformers, LLMs use a causal attention mask. As a result, the dense scheduling approach wastes computational resources on zero masks.

% Furthermore, the fusion of operations such as \( {\bf Q} \otimes {\bf k} \), softmax, and \( \bf S \otimes {\bf V} \) can reduce the additional accesses to DDR. The state-of-the-art kernel fusion implementation for resource-abundant GPUs is Flash Attention. However, GPU-optimized computation is not suitable for FPGAs, as GPUs have many more computational cores and much larger on-chip SRAM compared to the on-chip BRAM/URAM available on FPGAs. To address these challenges, we propose the reverse attention method, which utilizes fused attention and reverse reorder scheduling, specifically tailored for edge FPGAs.

% \begin{table}[ht]
% \caption{Comparison of different attention approaches.}
% \centering
% \resizebox{\linewidth}{!}{
% \begin{tabular}{|c|c|c|c|c|}
% \hline
% \textbf{Approach} & \textbf{Data Block Load} & \textbf{Iteration Count} & \textbf{Bandwidth} \\
% \hline
% Reverse Scheduling (Our Design)  & $\frac{n^2}{2p} + \frac{n}{2}  $ & $\frac{n^2}{2p} + \frac{n}{2}$ & $\sim 1$ \\
% \hline
% naive scheduling & $n^2 + n$ & $\frac{n^2}{p}$ & $\sim p$ \\
% \hline
% dense scheduling ~\cite{edgemoe}& $\frac{n^2}{p} + n + p - 1$ & $\frac{n^2}{p} + p - 1$ & $\sim 1$ \\

% \hline
% \end{tabular}
% }
% \label{tab:compare}
% \end{table}

% \begin{figure}[t] 
% \centering
% \includegraphics[width=0.85\linewidth]{Figures/dense_schedule.pdf}
% \caption{Dense attention scheduling ($p = 4$).}
% \label{fig:dense_schedule} 
% \end{figure}

% \begin{figure}[t] 
% \centering
% \includegraphics[width=\linewidth]{Figures/reverse_schedule.pdf}
% \caption{Reverse attention scheduling ($p = 4$).}
% \label{fig:reverse_schedule} 
% \vspace{1em}
% \end{figure}
% \subsubsection{Reverse Attention}
% % The pseudo-code of \textit{reverse attention} is shown in Algorithm \todo{algorithm} and the scheduling is shown in Fig. \ref{fig:reverse_schedule}. Assume that the current length of prefill tokens is $n$. $1< i \leq n$ and $1 < j \leq n$ indicate the current token index for $\bf q$ and $\bf k$, $\bf v$. There are $N_{h}$ heads in total.

% % The kernel fusion computation can be deemed a special case for Flash Attention V2 ~\cite{dao2022flashattentionfastmemoryefficientexact} when block size equals 1. The head-wise formula of the case of 2 consecutive block can be written as,  

% The reverse attention scheduling is depicted in Fig. \ref{fig:reverse_schedule}. Assume that the current length of the prefill tokens is \( n \), with \( 1 < i \leq n \) and \( 1 < j \leq n \) representing the current token indices for \( {\bf q} \) and \( {\bf k}, {\bf v} \), respectively. There are a total of \( N_{h} \) heads.

% The kernel fusion computation can be considered a special case of Flash Attention V2 ~\cite{dao2022flashattentionfastmemoryefficientexact} when the block size is equal to 1. The head-wise formula for the case with two consecutive blocks can be written as follows:

% \begin{equation}
% \left\{
% \begin{aligned}
% m^{(1)} &= s^{(1)} \\
% \ell^{(1)} &= e^{s^{(1)}-m^{(1)}} \\
%{ \bf o}^{(1)} &= e^{s^{(1)}-m^{(1)}} {\bf v}^{(1)} \\
% m^{(2)} &= \max \left(m^{(1)}, s^{(2)}\right) = m \\
% \ell^{(2)} &= e^{m^{(1)}-m^{(2)}} \ell^{(1)} + e^{s^{(2)}-m^{(2)}} \\
% &= e^{s^{(1)}-m} + e^{s^{(2)}-m} = \ell \\
% \bf{p}^{(2)} &= e^{s^{(2)}-m^{(2)}} / \ell^{(2)} \\
%{ \bf o}^{(2)} &={ \bf o}^{(1)} / e^{m^{(1)}-m^{(2)}} + e^{s^{(2)}-m^{(2)}} {\bf v}^{(2)} \\
% &= e^{s^{(1)}-m} {\bf v}^{(1)} + e^{s^{(2)}-m} {\bf v}^{(2)} \\
%{ \bf o}^{(2)} &={ \bf o}^{(2)} / \ell^{(2)} ={ \bf o}
% \end{aligned}
% \right.
% \tag{3}
% \end{equation}
% where \( m \) denotes the maximum value, \( s = {\bf q}_i \otimes {\bf k}_i \) represents the MAC result of the vector dot product, \( \ell \) is the denominator factor, and \({ \bf o} \) is the numerator vector. The upper index indicates the current step in the kernel fusion computation.

% Regarding scheduling, instead of starting from the first token \( {\bf q}_1 \), our schedule begins from \( {\bf q}_{n-1} \). Specifically, the level of parallelism is set to \( p \) (as illustrated in the figure for \( p = 4 \)). The factor \( p \) also implies that the on-chip BRAM can store \( p \) tokens of \( {\bf q}_i \). In each iteration, one \( {\bf q}_i \) token is loaded onto the on-chip memory (with the first batch loading \( {\bf q}_N \) to \( {\bf q}_{n-3} \)). Simultaneously, the corresponding \( {\bf k}_j \) and \( {\bf v}_j \) tokens are loaded for computation. After all \( n \) \( {\bf k}_j \) and \( {\bf v}_j \) tokens have been loaded and the fused-kernel computation is completed, the next iteration will evict \( p \) \( {\bf k}_j \) and \( {\bf v}_j \) tokens, starting from \( {\bf k}_{n-3} \) and \( {\bf v}_{n-3} \), to avoid redundant computations arising from the causal attention mask. The iteration continues until all \( 1 < i \leq n \), \( 1 < j \leq n \) are traversed. In this approach, the only required input buffers are for \( p \) \( {\bf q}_i \) tokens, one \( {\bf k}_j \), and one \( {\bf v}_j \). Additionally, the intermediate buffers include: \( N_{h} \times p \) multi-head MAC intermediate results \( s \), \( N_{h} \times p \) multi-head previous max values \( m \), and \( N_{h} \times p \) intermediate denominators \( \ell \).

% After introducing the reverse attention process, a comparison between naive attention, dense attention in Edge Moe ~\cite{edgemoe}, and our proposed reverse attention is provided in Table \ref{tab:compare}. The comparison indicates that the iteration count for reverse attention is the lowest, and the required bandwidth remains constant. Furthermore, the redundant outcome rate of the computation core can be directly assessed from the attention map in Fig. \ref{fig:attention_map}. The naive and dense scheduling approaches do not account for the causal mask, leading to many redundant computed values. 
% %It is important to note that due to kernel fusion, latency tests cannot be directly compared, and results will be presented in \todo{ablation study}.
% % where $m$ indicates the max value, $s = q_{i} \otimes k_{i}$ is the MAC result of vector dot product, $\ell$ is the denominator factor, $\bf o$ is the numerator vector. The upper index indicates the current step of kernel fusion computation.
% % As for scheduling, instead of starting from the first token $\bf q_{1}$, our schedule starts from the ${\bf q}_{n-1}$. In detail, the level of parallelism is set as $p$ (the figure illustrate the situation of $p = 4$). The $p$ factor also indicates that on-chip BRAM can stores $p$ $\bf q_{i}$ tokens. For each iteration, one $\bf q_{i}$ token is loaded to the onchip (the first batch will be ${\bf q}_{n}$ to ${\bf q}_{n-3}$). Meanwhile, for every iteration, the corresponding ${\bf k}_{j}$ and ${\bf v}_{j}$ are loaded for computation. After all $n$ ${\bf k}_{j}$ ${\bf v}_{j}$ tokens are loaded and finished fused-kernel computation, next iteration will evict p ${\bf k}_{j}$ ${\bf v}_{j}$ and start from ${\bf k}_{n - 3}$ ${\bf v}_{n-3}$ to avoid redundant computation that stems from the casual attention mask. The iteration carries on until all $1 < i \leq n$, $1 < j \leq n$ are traversed. In this way, the only required input buffers are for $p$ ${\bf q}_{i}$, one ${\bf k}_{j}$ and one ${\bf v}_{j}$. Moreover, the intermediate buffers are: $N_{h} \times p$ multi-head MAC intermediate results $s$, $N_{h} \times p$ multi-head previous max value $m$, and $N_{h} \times p$ intermediate denominator $\ell$.

% % After introducing the process of reverse attention, a comparison between naive attention, dense attention in Edge Moe ~\cite{edgemoe} and our proposed reverse attention can be given in \ref{tab:compare}. The comparison indicates that the iteration count for reverse attention is the lowest, and the required bandwidth is constant. Moreover, the redundant outcome rate of the computation core can be directly compared from the attention map in Fig. \ref{fig:attention_map}. The naive and dense scheduling didn't consider the causal mask and many computed values are redundant. Notice that due to the kernel fusion, latency test can not be directly compared and will be displayed in \todo{ablation study}. 
% % Table \todo{Bandwidth/IDLE RATE analysis} displays the general case

% %\subsection{Unified Decoding Attention and LM-Head}
% \subsection{Hardware Specialization and Reuse for Decoding-Phase Attention and LM Head}

% The previous section focuses on optimizing the Prefill phase in LLM inference. Since we consider efficient end-to-end LLM inference on single edge Device, both Prefill and Decoding phases require efficient implementation to maximize global performance.
% % in this section, we optimize the Decoding phase, another significant part of inference. The attention mechanism %(Attention(Q,k,V)=softmax(dk QKT)V) 
% The attention mechanism is a core component in both phases, but its computational characteristic differs. In the prefill phase, calculating attention involves operations on matrices representing the entire input sequence. In the decoding phase, it involves operations between the vector representation of the single new token and the cached matrices of keys and values. This difference presents an opportunity for hardware specialization.

% In the decoding phase, a single new token is generated per step. Let the total sequence length (prompt + already generated tokens) be $m$. The query $\bf q$ is now a $1 \times n$ vector corresponding to the new token. The key ($K_{cache}$) and value ($V_{cache}$) matrices contain the cached representations of all m previous tokens and are of dimensions m×d. The core attention computation involves:

% 1.	Calculating attention scores: ${\bf q} \times K_{cache}^T$ (a $1 \times n$ vector multiplied by a $n\times m$ matrix, resulting in a ${1\times m}$ score vector).
% 2.	Applying softmax to the scores.
% Multiplying the resulting $1\times m$ vector by $V_{cache}$ (an $m\times n$ matrix) to get the $1\times n$ output vector.
% The computation involves primarily matrix-vector and vector-vector operations. The computational load per step ($O(Nd)$) is significantly lower than the total prefill computation. However, this phase requires fetching the large $K_{cache}$ and $V_{cache}$ matrices from memory (e.g., off-chip DDR) in every step. Consequently, the decoding phase is often memory-bandwidth bound, especially as the sequence length $m$ grows. 
% Since computation is less intensive and latency is dominated by memory access (fetching the KV cache), massive parallelism is inefficient and wastes resources. A more sequential or lower-parallelism computation unit is sufficient. This approach significantly reduces the required on-chip hardware resources (e.g., number of PEs, buffer sizes) compared to the prefill unit. 
% The profiling result in Fig.~\ref{fig:load_compute} clearly shows that the attention in the decoding phase is memory-bounded, while the prefill phase is compute-bounded. This demonstrates the necessity of lightweight decoding implementation to save on-chip resources for the Prefill phase.


% \begin{figure}[ht]
%     \centering

%     \includegraphics[width=\linewidth]{Figures/llm_load_compute_comparison_font16.png}

%     %\caption{The Structured Data Access for SpMV}
%     \caption{Characterization of Attention Module during Prefill/Decoding Phase }
%     %%\Description{. }
%     \label{fig:load_compute}
% \end{figure}


% \subsubsection{Reuse of Decoding Attention for LM Head}
% Following the processing 
% thr-ough n transformer blocks in typical LLM architectures like LLaMA, the final step before token generation involves the LM Head. This component performs a crucial linear projection, mapping the final hidden state output from the last transformer block to a vector of logits representing the probability distribution over the entire vocabulary.
% For inference, particularly during the auto-regressive decoding phase, where one token is generated at a time, the input to the LM Head is the final hidden state vector corresponding to the token being predicted. This vector has dimensions $[1,n]$, where n represents the hidden state dimension (e.g., $n=1536$ for some models). The LM Head uses a large weight matrix of dimensions $[n,V]$, where $V$ is the vocabulary size (e.g., $V=32000$), to compute the output logit vector of dimensions $[1,V]$. The core computation is thus a matrix-vector multiplication: $[1,n]\times[n,V]\rightarrow [1,V]$.
%  This computation (O(HV)) has similar characteristics to decoding attention: it's a matrix-vector operation with limited data reuse opportunities, heavily reliant on fetching a large matrix (the $K_{cache}$ or the LM Head weights), making it memory-bound (especially as V $\gg$ H).
 
% Given this similarity, we reuse the decoding-phase attention hardware to execute the LM Head computation. This eliminates the need for a dedicated LM Head unit, yielding substantial area and power savings. The performance impact is negligible because the LM Head executes only once per generated token, whereas attention occurs n times (once per layer). 


% Reusing the Decoding-phase Attention hardware for both attention and the LM Head precludes a fully fused attention pipeline within it. We therefore adopt a decoupled execution model for attention sub-steps during decoding:
% (1) Attention Score Computation: ${\bf s}={\bf q}\times K_{cache}^T$ .
% (2) Softmax: ${\bf p}=softmax({\bf s})$.
% (3) Value Aggregation: Compute the final attention output ${\bf o}={\bf p}\times V_{cache}$.
% This is efficient because the intermediate score/probability vector (1×m) is small enough to be buffered on-chip BRAM with minimal latency penalty.
% On the contrary, for the Prefill phase, the intermediate n×n attention score matrix would be far too large for practical on-chip buffering. Thus it requires a fully fused attention pipeline that integrates attention score calculation, softmax, and value aggregation in a single, uninterrupted hardware pass to minimize off-chip data movement.

% % For the softmax step, we employ a numerically stable online softmax approach to prevent potential overflow/underflow issues common with standard softmax calculation on hardware, ensuring robust computation.

% % \begin{algorithm}[N_{h}]
% % \For{$i = 1$ to $m$}
% %     \State $e_{shifted}[i] \leftarrow \exp(x_i - m)$ \Comment{Compute exponent of shifted value}
% %     \State $S \leftarrow S + e_{shifted}[i]$ \Comment{Update sum}
% % \EndFor
% % \end{algorithm}



% % Following the processing 
% % thr-ough n transformer blocks in typical LLM architectures like LLaMA, the final step before token generation involves the LM Head. This component performs a crucial linear projection, mapping the final hidden state output from the last transformer block to a vector of logits representing the probability distribution over the entire vocabulary.
% % For inference, particularly during the auto-regressive decoding phase, where one token is generated at a time, the input to the LM Head is the final hidden state vector corresponding to the token being predicted. This vector has dimensions $[1,n]$, where n represents the hidden state dimension (e.g., $n=1536$ for some models). The LM Head uses a large weight matrix of dimensions $[n,V]$, where $V$ is the vocabulary size (e.g., $V=32000$), to compute the output logit vector of dimensions $[1,V]$. The core computation is thus a matrix-vector multiplication: $[1,n]\times[n,V]\rightarrow [1,V]$.
% % We observe that the fundamental computational pattern of the LM Head during decoding – a matrix-vector multiplication – precisely mirrors the nature of the attention score calculation (${\bf q}\times\text{Kcache}^T$) within the decoding phase, as handled by our proposed Decoding Optimization Unit (DOU). Both operations involve multiplying a single vector (the query $\bf q$ or the hidden state) by a large matrix (the Kcache or the LM Head weights). Furthermore, both exhibit similar computational characteristics: they offer limited data reuse opportunities and are significantly constrained by memory bandwidth, primarily due to the need to load the large weight matrices (Kcache or the $[n,V]$ LM Head matrix) from memory. Notably, the LM Head's weight matrix dimension $V$ is typically much larger than the hidden dimension $n$, exacerbating its memory-bound nature.
% % Given these strong similarities in computational structure and characteristics, we propose to reuse the Decoding Attention Unit, originally designed for efficient decoding-phase attention, to also execute the LM Head computation. This hardware reuse strategy offers significant advantages:

% % 1.	Resource Efficiency: It eliminates the need for a separate, dedicated hardware unit for the LM Head projection. This directly translates to substantial savings in on-chip silicon area and associated static power consumption, contributing to a more compact and efficient accelerator design.

% % 2.	Minimal Performance Impact: The potential performance overhead introduced by this reuse is negligible. While the attention mechanism is computed n times (once per transformer layer) for each generated token, the LM Head computation is performed only once per token, at the very end of the forward pass. Therefore, scheduling the LM Head calculation on the the Decoding Attention Unit introduces minimal contention and has an insignificant impact on the overall token generation latency compared to the cumulative latency of the n transformer blocks.

% % By leveraging the the Decoding Attention Unit, which is already optimized for the memory-bound matrix-vector operations characteristic of the decoding phase, we can efficiently execute the LM Head projection without dedicating additional specialized hardware. This approach aligns with our overall design philosophy of maximizing hardware utilization and minimizing resource redundancy by exploiting computational parallels within the LLM inference workflow.

% % \subsubsection{Online Softmax}
% % The process of attention computation can be broken down into sequential stages:
% % (1) Attention Score Computation: ${\bf s}={\bf q}\times K_{cache}^T$ .
% % (2) Softmax: ${\bf p}=softmax({\bf s})$.
% % (3) Value Aggregation: Compute the final attention output ${\bf o}={\bf p}\times V_{cache}$.

% % For the Prefill phase, the intermediate n×n attention score matrix would be far too large for practical on-chip buffering. Therefore, we use a fully fused attention pipeline that integrates attention score calculation, softmax, and value aggregation ($\bf p\times V_{cache}$) in a single, uninterrupted hardware pass to minimize off-chip data movement.
% % However, for the decoding phase, the intermediate data transferred between these stages is the attention score vector s (or the resulting probability vector {\bf p}), which has dimensions $[1 \times m]$ (where m is the current sequence length). This single vector is compact enough to be stored effectively in on-chip memory resources, such as Block RAM (BRAM) or registers, with negligible latency overhead associated with the temporary storage and retrieval. Consequently, we adopt a decoupled execution model for the sub-steps of the attention mechanism during the decoding phase. It enables the modular hardware reuse for the decoding phase attention and LM Head, maximizing resource utilization with minimal impact on performance due to efficient on-chip handling of the intermediate vector.





% % The strategic decision to reuse the Decoding Attention Unit hardware for both decoding-phase attention and the subsequent LM Head computation influences the optimal dataflow for the attention mechanism itself. Specifically, this reuse prevents the implementation of a fully fused attention pipeline within the DOU – one that would seamlessly integrate attention score calculation, softmax, and value aggregation ($\bf p\times V_{cache}$) in a single, uninterrupted hardware pass. Such fusion is precluded because the unit must be available for the distinct LM Head task.

% % Consequently, we adopt a decoupled execution model for the sub-steps of the attention mechanism during the decoding phase when utilizing the DOU. The process is broken down into sequential stages:

% % Attention Score Computation: Calculate the raw attention scores ${\bf s}={\bf q}\times K_{cache}^T$ .
% % Softmax Normalization: Apply the softmax function to the score vector s to obtain probabilities ${\bf p}=softmax({\bf s})$.
% % Value Aggregation: Compute the final attention output vector ${\bf o}={\bf p}\times V_{cache}$.
% % This decoupling is highly practical and efficient specifically within the decoding context. The intermediate data transferred between these stages is the attention score vector s (or the resulting probability vector {\bf p}), which has dimensions $[1 \times m]$ (where m is the current sequence length). This single vector is compact enough to be stored effectively in on-chip memory resources, such as Block RAM (BRAM) or registers, with negligible latency overhead associated with the temporary storage and retrieval. This feasibility contrasts sharply with the prefill phase, where the intermediate n×n attention score matrix would be far too large for practical on-chip buffering between decoupled steps, making fusion more critical there. The primary benefit of this decoupled approach in the decoding phase is that it enables the modular hardware reuse for the DOU and LM Head, maximizing resource utilization with minimal impact on performance due to efficient on-chip handling of the intermediate vector.

% % For the critical Softmax Normalization stage within this decoupled flow, we implement an online softmax algorithm. The standard softmax formulation, $p_i =e^{x_i} /sum(\sum_{1}^{j}e^{x_j})$ , can suffer from numerical instability if the input scores ($x_i$) have large magnitudes, potentially causing overflow during the exponentiation or underflow/loss of precision during the summation if scores are large and negative. The online approach enhances numerical stability by iteratively processing the input score vector and maintaining running statistics. It typically involves normalizing values relative to the running maximum score encountered so far. This avoids the computation of potentially very large exponential values directly, leading to a more robust hardware implementation.








% \subsection{Implementation and Optimization of Special Function Units}

% Besides the core modules discussed in previous sections, LLM inference also relies on several essential special functions, including Quantization / Dequantization, RMSNorm, and Activation Function. 
% % While computationally less intensive than the core modules and representing a smaller fraction of total inference time, their efficient implementation is crucial for overall performance and resource utilization. 
% These operations are computationally less intensive, therefore, our strategy focuses on lightweight hardware implementations and operator fusion to minimize overhead. For efficient data handling, vector data for these modules is processed in 256-bit packets, aligning with AXI bus width.

% \textbf{Quantization \& Dequantization}: 
% % Our use of ternary Linear modules requires quantization on activations. We employ INT8 symmetric absolute maximum (ABSMAX) quantization. 
% The activations need to be quantized before the ternary Linear modules. We employ ABSMAX Quantization, which involves two passes: (1) finding the maximum absolute value to compute the scale factor, and (2) applying this scale element-wise. Fig.~\ref{Fig:Bitnet} shows that quantization follows the RMSNorm. We fuse these operations to reduce data movement. The dequantization is fused into the Linear output pipeline.

% \textbf{RMSNorm}:
% RMSNorm involves two passes. The first calculates the Root Mean Square of the input x: $RMS(x) =\sqrt{\frac{1}{n}\sum_{i =1}^{n}x_{i}^{2}}$. The second pass normalizes the input by dividing by the RMS value and multiplying by a learned scaling parameter $\gamma$. Recognizing that both RMSNorm and ABSMAX Quantization involve a two-pass traverse, we fuse these four logical steps into two optimized hardware passes. This significantly minimizes the data movement.

% \textbf{Activation function}:
% % The SILU activation function, defined as: SILU(x)=x⋅σ(x)=x⋅1+e−x1 is applied element-wise, typically to the output of the gate projection within the Feed-Forward Network (FFN) block. We implement SILU as a pipelined element-wise unit. To hide its latency, we fuse this computation directly into the output pipeline of the preceding Linear module responsible for generating the gate projection tensor. This allows the SILU computation to overlap with the final stages of the Linear operation, effectively minimizing its impact on the critical path.
% % The element-wise SILU activation (SILU(x)=x⋅σ(x)) is pipelined and fused directly into the output stage of the generating Linear module, effectively hiding its latency.
% The element-wise SILU activation ($x\cdot\frac{1}{1+e^{-x}}$) is required after the Gate projection in Feed-Forward Network (FFN) block. The SILU is pipelined and fused directly into the preceding Linear module, effectively hiding its latency.


\begin{table*}[t]
\centering
\caption{Unified cross-platform and FPGA-based comparison for edge LLM inference. Resource utilization is reported for FPGA works. Throughput (TK/S) is tokens per second; energy efficiency (TK/J) is tokens per joule. Max DDR bandwidths are theoretical.}
\vspace{-2mm}
\label{tab:unified-llm-comparison}
\setlength{\tabcolsep}{3pt}
\renewcommand{\arraystretch}{0.85}
\resizebox{\textwidth}{!}{%
\begin{tabular}{c c c c c c c c c c c c c c c c}
\toprule
\multirow{2}{*}{\textbf{Work}} &
\multirow{2}{*}{\textbf{Platform}} &
\multirow{2}{*}{\textbf{Processor}} &
\multirow{2}{*}{\textbf{Model}} &
\multirow{2}{*}{\textbf{Precision}} &
\multirow{2}{*}{\textbf{Max DDR Bandwidth(GB/s)}} &
\multicolumn{5}{c}{\textbf{FPGA Resource Utilization}} &
\multirow{2}{*}{\textbf{Power (W)}} &
\multicolumn{2}{c}{\textbf{Throughput (TK/S)}} &
\multicolumn{2}{c}{\textbf{Energy Efficiency (TK/J)}} \\
\cmidrule(lr){7-11} \cmidrule(lr){13-14} \cmidrule(lr){15-16}
& & & & & & \textbf{LUT} & \textbf{FF} & \textbf{DSP} & \textbf{BRAM} & \textbf{URAM} & & \textbf{Prefill} & \textbf{Decode} & \textbf{Prefill} & \textbf{Decode} \\
\midrule
Raspberry Pi 5 \cite{adafruit2025qwen3} & SoC & $4\times$ \ Cortex-A76 & Qwen 0.6B & W4-A16 & 17.1 & — & — & — & — & — & 7.8 & 61.8 & 16.6 & 7.92 & 2.12 \\
Jetson Orin Nano \cite{jetsonailab2025slm} & GPU SoC & $8\times$ \ GPU SM & TinyLLaMA 1.1B & W4-A16 & 68.3 & — & — & — & — & — & 25 & 324.9 & 67.6 & 12.9 & 2.70 \\
\midrule
SECDA \cite{haris2024designing} & FPGA SoC & $2\times$ \ Cortex-A53 & TinyLLaMA 1.1B & W4-A16 & 2.1 & — & — & — & — & — & — & — & 0.6 & — & — \\
LLaMAF \cite{LLaMAf} & FPGA SoC & ZCU102 & TinyLLaMA 1.1B & W8-A8 & 19.2 & 164K & 171K & 528 & 223 & — & 5.1 & — & 1.5 & — & 0.29 \\
MEADOW \cite{moitra2025meadow} & FPGA SoC& ZCU102 & OPT 1.3B & W8-A8 & 19.2 & 150K & — & 845 & 2034 & — & 10 & 100 & 2 & 10 & 0.20 \\
\midrule
\textbf{TeLLMe (KV260)} & FPGA SoC & — & BitNet 0.73B & W1.58-A8 & 17.1 & 98K & 137K & 610 & 98.5 & 60 & \textbf{4.8} & \textbf{143} & \textbf{25} & \textbf{29.8} & \textbf{5.2} \\
\bottomrule
\end{tabular}
}
% \begin{tablenotes}
% \tiny
% \item[*] *MEADOW implements the INT8 Matmul Unit and no for profiling.
% \end{tablenotes}
\end{table*}

% \begin{table*}[t]
% \centering
% \begin{threeparttable}
% \caption{Unified cross-platform and FPGA-based comparison for edge LLM inference. Resource utilization is reported for FPGA works. Throughput (TK/S) is tokens per second; energy efficiency (TK/J) is tokens per joule. Max DDR bandwidths are theoretical.}
% \label{tab:unified-llm-comparison}
% \setlength{\tabcolsep}{3pt}
% \renewcommand{\arraystretch}{0.85}
% \resizebox{\textwidth}{!}{%
% \begin{tabular}{c c c c c c c c c c c c c c c c}
% \toprule
% \multirow{2}{*}{\textbf{Work}} &
% \multirow{2}{*}{\textbf{Platform}} &
% \multirow{2}{*}{\textbf{Processor}} &
% \multirow{2}{*}{\textbf{Model}} &
% \multirow{2}{*}{\textbf{Precision}} &
% \multirow{2}{*}{\textbf{Max DDR Bandwidth(GB/s)}} &
% \multicolumn{5}{c}{\textbf{FPGA Resource Utilization}} &
% \multirow{2}{*}{\textbf{Power (W)}} &
% \multicolumn{2}{c}{\textbf{Throughput (TK/S)}} &
% \multicolumn{2}{c}{\textbf{Energy Efficiency (TK/J)}} \\
% \cmidrule(lr){7-11} \cmidrule(lr){13-14} \cmidrule(lr){15-16}
% & & & & & & \textbf{LUT} & \textbf{FF} & \textbf{DSP} & \textbf{BRAM} & \textbf{URAM} & & \textbf{Prefill} & \textbf{Decode} & \textbf{Prefill} & \textbf{Decode} \\
% \midrule
% Raspberry Pi 5 & SoC & $4\times$ Cortex-A76 & Qwen 0.6B & W4-A16 & 17.1 & — & — & — & — & — & 7.8 & 61.8 & 16.6 & 7.92 & 2.12 \\
% Jetson Orin Nano & GPU SoC & $8\times$ GPU SM & TinyLLaMA 1.1B & W4-A16 & 68.3 & — & — & — & — & — & 25 & 324.9 & 67.6 & 12.9 & 2.70 \\
% \midrule
% SECDA \cite{haris2024designing} & FPGA SoC & $2\times$ Cortex-A53 & TinyLLaMA 1.1B & W4-A16 & 2.1 & — & — & — & — & — & — & — & 0.6 & — & — \\
% LLaMAF \cite{LLaMAf} & FPGA SoC & ZCU102 & TinyLLaMA 1.1B & W8-A8 & 19.2 & 164K & 171K & 528 & 223 & — & 5.1 & — & 1.5 & — & 0.29 \\
% MEADOW \tnote{a} \cite{moitra2025meadow} & FPGA SoC & ZCU102 & OPT 1.3B & W8-A8 & 19.2 & 150K & — & 845 & 2034 & — & 10 & 100 & 2 & 10 & 0.20 \\
% \midrule
% \textbf{TeLLMe (KV260)} & FPGA SoC & — & BitNet 0.73B & W1.58-A8 & 17.1 & 98K & 136K & 610 & 98.5 & 60 & \textbf{4.8} & \textbf{143} & \textbf{25} & \textbf{29.8} & \textbf{5.2} \\
% \bottomrule
% \end{tabular}
% }
% \begin{tablenotes}
% \small
% \item[a] This is a footnote for item MEADOW.
% \end{tablenotes}
% \end{threeparttable}
% \end{table*}


\begin{table}[t]
\centering
\caption{Perplexity and Intelligence/J comparison on WikiText-2. Intelligence/J is defined as $\frac{tokens/s}{(PPL \cdot Power)}$~\cite{tenent}. Throughput (TK/S) is tokens per second.}
\vspace{-2mm}
\label{tab:ppl-intelligenceJ}
\setlength{\tabcolsep}{4pt}
\renewcommand{\arraystretch}{0.85}
\resizebox{0.48\textwidth}{!}{%
\begin{tabular}{c c c c c c c c}
\toprule
\multirow{2}{*}{\textbf{Work \& Platform}} &
\multirow{2}{*}{\textbf{Model}} &
\multirow{2}{*}{\textbf{Power (W)}} &
\multirow{2}{*}{\textbf{WT-2 (PPL)}} &
\multicolumn{2}{c}{\textbf{Throughput (TK/S)}} &
\multicolumn{2}{c}{\textbf{Intelligence/J}} \\
\cmidrule(lr){5-6} \cmidrule(lr){7-8}
& & & & \textbf{Prefill} & \textbf{Decode} & \textbf{Prefill} & \textbf{Decode} \\
\midrule
Raspberry Pi 5 & Qwen 0.6B & 7.8  & 24.00  & 61.8   & 16.6  & 0.330 & 0.089 \\
Jetson Orin Nano & TinyLLaMA 1.1B & 25.0 & 12.42  & 324.9  & 67.6  & 1.046 & 0.218 \\
SECDA (Pynq) & TinyLLaMA 1.1B & 1.2    & 12.42  & —      & 0.6   & —     & 0.040  \\
LLaMAF (ZCU102) & TinyLLaMA 1.1B & 5.1    & 8.89 & —      & 1.5   & —     & 0.041     \\
MEADOW (ZCU102) & OPT 1.3B & 10.0 & 15.41  & 100.0  & 2.0   & 0.649 & 0.013 \\
\textbf{TeLLMe (KV260)} & \textbf{BitNet 0.73B} & \textbf{4.8} & \textbf{12.79} & \textbf{143} & \textbf{25} & \textbf{2.330} & \textbf{0.407} \\
\bottomrule
\end{tabular}
 }
\end{table}
%\vspace{-1mm}
\section{Experimental Results and Analysis} \label{expr}
%\vspace{-1mm}
\subsection{Experiment Setup}
%\vspace{-1mm}
% \todo[inline]{give the hardware evaluation setups}
We implement the TeLLMe accelerator using high-level synthesis in C/C++ with Vitis HLS 2024.1 and Vivado 2024.1. The design is evaluated on the AMD Kria KV260 platform (Zynq UltraScale+ XCK26 MPSoC). To achieve timing closure, a clock frequency of 250 MHz is employed for the final bitstream generation. The end-to-end latency results are measured with the PYNQ Runtime library, while the latency breakdown is derived from cycle-accurate RTL simulations. For fair comparison, we select prior works that implement approximately 1B-parameter models on edge platforms.
The implemented model is BitNet 0.73B~\cite{bitnet158}, comprising 49M parameters for the language model head weights and embedding token table, alongside 680M parameters for the Transformer decoder layer weights.
According to the parameter selection formulas in Sec.~\ref{SEC:PARAMETER_WBMU}, the parameters for our harderware accelerator are $G = 3$, $T = 28$, $Q = 16$, $U = 48$, $B_{\text{idx}} = 5$, with the weight index vector ${\bf w}_{\text{idx}}$ having a 140-bit width.
Limited by routing and placement constraints for an extra module, we offload the language model head (LM Head) computation to the Processing System (PS) side. Using ARM NEON SIMD computing kernel with W8A8 quantization, the LM Head latency is 9 ms, with negligible accuracy loss, as shown in Table~\ref{tab:ppl-intelligenceJ}. This latency is considered in our end-to-end performance results.

\begin{figure}[h]
    \centering
    \includegraphics[width=0.95\linewidth]{Figures/TELLME_device_map_NEW_v2.pdf}
    %\caption{The Structured Data Access for SpMV}
    \caption{TeLLMe Hardware Accelerator Floorplan on KV260}
    %%\Description{}
    \label{fig:device_map}
\end{figure}


% We implement the TeLLMe accelerator using high-level synthesis C/C++ in Vitis HLS 2024.1  and Vivado 2024.1. We evaluate our design on Kria KV260 (Zynq UltraScale+ XCK26 MPSoC). To ensure timing closure, we use 250 MHz for the final bitstream generation. The latency results are measured through the Pynq platform, and the latency breakdown comes from the RTL simulation. To ensure fair comparsion, we choose works that implement around 1B model on edge platforms.

% The implemented model is BitNet 0.73B, with 49M language model head weight and embedding token table, as well as 680M decoder layer weight. 
% According the parameter selection formulas in Sec. \ref{SEC:PARAMETER_WBMU}, the parameter of our model is $G = 3$, $T = 28$, $Q = 16$, $U = 48$, $B_{\text{idx}} = 5$, the weight index vector ${\bf w}_{\text{idx}}$ is with 140 bitwidth.
\subsection{TeLLMe Inference Performance}
%\vspace{-1mm}
Fig.~\ref{fig:perf} presents the performance characteristics of our TeLLMe accelerator across ten different [prompt, generation] configurations. The results demonstrate a clear inverse relationship between sequence length and performance metrics. Our design achieves up to 25 tokens/s throughput with great real word prefill latency. 

Notably, configurations with prompt lengths less than 256 tokens achieve decoding throughput above 16 tokens/s and prefill latency below 2.25s, demonstrating TeLLMe's performance under limited resource \& power budget. These results also suggest that practical FPGA deployments should target moderate sequence lengths to balance latency requirements with inference throughput, achieving optimal efficiency ratios exceeding 20 tokens/s per prefill second for configurations [64,128] through [128,256].

Our inference latency breakdown in Fig.~\ref{fig:breakdown} reveals distinct computational characteristics between the prefill and decoding phases on FPGA. The decoding phase is predominantly memory-bound, with performance limited by memory bandwidth due to weight loading and KV cache accesses, while the prefill phase is compute-bound, requiring substantial parallel matrix operations. By leveraging the KV cache, decoding attention processes only a single new token against cached keys and values, consuming significantly less time per token than prefill, which recomputes attention across the entire sequence. This fundamental difference motivates our architecture design.
The further resource breakdown is shown in Table \ref{tab:utilization}. 

% Regarding BRAM usage, most of it is consumed by the top-level AXI buffer. DSP resources are mainly utilized by the Attention modules due to their FP16 precision. LUTs are primarily consumed by the TL-based matrix multiplication unit for the TL tables. URAM is used by the matrix multiplication weight buffer to support ping-pong operation.


\begin{figure}[h]
    \centering

    \includegraphics[width=\linewidth]{Figures/llm_performance_chart_Ye.pdf}

    %\caption{The Structured Data Access for SpMV}
    \caption{TeLLMe LLM Inference Performance}
    %%\Description{}
    \label{fig:perf}
\end{figure}

\begin{figure}[ht]
    \centering

    \includegraphics[width=\linewidth]{Figures/latency-breakdown-v2.pdf}

    %\caption{The Structured Data Access for SpMV}
    \caption{TeLLMe LLM Inference Latency Breakdown (Length = 128)}
    %%\Description{}
    \label{fig:breakdown}
\end{figure}
% \begin{table}[N_{h}]
% \centering
% \caption{test}
% \resizebox{\columnwidth}{!}{%
% \begin{tabular}{|c|c|c|c|c|c|}
% \hline
% Module                     & BRAM\_18K  & DSP        & FF            & LUT           & URAM      \\ \hline
% Control \& Data Transfer   & 120        & 0          & 24973         & 5897          & 0         \\ \hline
% Attention (Prefill)  & 46         & 122        & 25629         & 35069         &           \\ \hline
% Attention (Decoding) & 24         & 134        & 17465         & 7028          & 0         \\ \hline
% Lookup-based Linear Unit   & 0          & 0          & 35765         & 54094         & 48        \\ \hline
% RMSNorm                    & 16         & 28         & 6202          & 5933          &           \\ \hline
% Misc (Add, Mul, RoPE, etc) & 0          & 72         & 45804         & 4973          & 0         \\ \hline
% Total                      & 206 (71\%) & 356 (28\%) & 155838 (66\%) & 112994 (93\%) & 48 (75\%) \\ \hline

% \end{tabular}%
% }
% \label{tab:utilization}
% \end{table}





% Please add the following required packages to your document preamble:
% \usepackage{graphicx}
\begin{table}[h]
\centering
\caption{FPGA resource consumption breakdown. Percentages at bottom indicate design utilization per resource.}
\label{tab:utilization}
\vspace{-2mm}
\setlength{\tabcolsep}{5pt}
\renewcommand{\arraystretch}{0.8}
\resizebox{\columnwidth}{!}{%
\begin{tabular}{lrrrrr}
\toprule
\textbf{Module} & \textbf{LUT} & \textbf{FF} & \textbf{BRAM} & \textbf{URAM} & \textbf{DSP} \\
\midrule
TLMM\texttt{-}FUSE Unit          & 43{,}137 & 51{,}894 &  5.5 &  0  & 320 \\
RPA Unit (Prefill Attention)        & 18{,}072 & 31{,}212 & 43.5 &  4  & 115 \\
DA Unit (Decoding Attention)       &  9{,}570 & 20{,}138 & 11.5 &  4  & 123 \\
RMS-MAX Unit                          &  6{,}092 & 11{,}175 &  4.0 &  4  &  47 \\
WBMU Unit                   &  6{,}885   &  1{,}614  & —    &  48   &  — \\
Control \& Communication                  & 9{,}114 & 13{,}412 & 34.0 & —  &   5 \\
Misc                             &  5{,}433 &  7{,}276 &   —  &  —  &   — \\
\midrule
\textbf{Total}                   & \textbf{98{,}303} & \textbf{136{,}721} & \textbf{98.5} & \textbf{60} & \textbf{610} \\
\textit{Utilization}             & \textit{(84\%)}   & \textit{(28\%)}     & \textit{(68\%)} & \textit{(94\%)} & \textit{(49\%)} \\
\bottomrule
\end{tabular}%
}
\vspace{-2mm}
\end{table}
\begin{table}[h]
\caption{LUT usage breakdown for different TLMM Designs}
\vspace{-2mm}
\label{tab:LUT}
\vspace{-2mm}
\centering
\setlength{\tabcolsep}{5pt}
\renewcommand{\arraystretch}{0.95}
\resizebox{\columnwidth}{!}{%
\begin{tabular}{lrrrr}
\toprule
\textbf{Design} & $\textbf{LUT}_{\text{PRE}}$ & $\textbf{LUT}_{\text{TL}}$ & $\textbf{LUT}_{\text{LPL}}$ & \textbf{TOTAL} \\
\midrule
None Table Lookup (Method 1) & -- & -- &  -- & 43,176 \\
Partial Storage Table Lookup (Method 2)& 3,117 & 6,440 & 25,643 & 35,200 \\
Full Storage Table Lookup (Ours) & 5,301  & 11,452 & 6,329 & 23,082 \\
\bottomrule
\end{tabular}
}
% \vspace{-5mm}
\end{table}

\subsection{Cross-Platform and FPGA-based Comparison for Edge LLM Inference}
%\vspace{-1mm}
We present a comprehensive evaluation of TeLLMe against state-of-the-art edge LLM accelerators across FPGA, CPU, and GPU platforms. Table~\ref{tab:unified-llm-comparison} summarizes the hardware characteristics and performance metrics, while Table~\ref{tab:ppl-intelligenceJ} provides quality-aware efficiency comparisons using the Intelligence/J metric, defined as throughput normalized by perplexity and power consumption.

\subsubsection{Comparison with FPGA-based Solutions}

TeLLMe demonstrates substantial advantages over existing FPGA implementations in both throughput and energy efficiency. Compared to SECDA, which achieves 0.6 tokens/s decoding throughput on a Pynq-Z2 board, TeLLMe delivers 41.7$\times$ higher decoding throughput (25 tokens/s) while maintaining comparable power consumption (4.8W). This significant performance gain stems from our optimized BitNet architecture and efficient memory subsystem design that better exploits the available memory bandwidth.

When compared to LLaMAF on the higher-end ZCU102 platform, TeLLMe achieves 16.7$\times$ higher decoding throughput despite running on the more resource-constrained KV260 SoC. Although LLaMAF benefits from the ZCU102's superior memory bandwidth (19.2 GB/s vs 17.1 GB/s) and operates at slightly higher power (5.1W), TeLLMe's energy efficiency is 17.9$\times$ better during decoding operations. This improvement is primarily attributed to BitNet's 1.58-bit weight quantization, which dramatically reduces memory traffic compared to conventional 8-bit quantization.

MEADOW implements a well-designed INT8 matmul and attention engine. However, it lacks detailed descriptions of FP operations, including quantization and dequantization processes, to ensure smooth quantization support. Meanwhile, it represents a more computation-intensive approach, deploying a larger OPT 1.3B model with extensive DSP utilization (845 DSPs) despite delivering worse perform model (15.41 vs TeLLMe 12.79 of WikiText-2 PPL). While MEADOW achieves fair prefill throughput (100 tokens/s), its decoding performance (2 tokens/s) is 12.5$\times$ lower than TeLLMe. More critically, TeLLMe achieves 26$\times$ higher decoding energy efficiency (5.2 TK/J vs 0.20 TK/J), demonstrating the effectiveness of our memory-optimized architecture for the memory-bound decoding phase. The Intelligence/J metric further highlights this advantage: TeLLMe achieves 31.3$\times$ better decoding intelligence efficiency (0.407 vs 0.013), indicating superior quality-normalized energy efficiency. 

From a resource utilization perspective, TeLLMe achieves these performance gains while maintaining modest FPGA resource consumption. Our design uses 98K LUTs and 610 DSPs, comparable to or lower than existing solutions, while leveraging URAMs (60 blocks) for efficient on-chip weight storage. This balanced resource allocation enables deployment on mid-range FPGA SoCs without requiring high-end devices. 

It is worth noting that both LLaMAF and MEADOW focus exclusively on accelerating the decoding stage, assuming that prefill operations will be offloaded to separate devices such as host CPUs. In contrast, TeLLMe provides a unified acceleration framework that efficiently handles both prefill and decoding phases on a single FPGA device. This holistic approach eliminates the need for heterogeneous processing pipelines and associated data transfer overhead. Despite accelerating both inference stages, TeLLMe uses fewer hardware resources while delivering superior throughput in both phases, achieving 143 tokens/s prefill and 25 tokens/s decoding performance. This comprehensive acceleration capability, combined with leading energy efficiency across all metrics, establishes TeLLMe as the state-of-the-art solution for FPGA-based edge LLM inference.

\subsubsection{Comparison with CPU and GPU Edge Platforms}
While GPU and CPU platforms offer higher absolute throughput due to greater computational resources and memory bandwidth, TeLLMe provides compelling advantages in energy efficiency and cost-effectiveness. The Jetson Orin Nano, equipped with 8 GPU streaming multiprocessors (SM) and 68.3 GB/s memory bandwidth, achieves 67.6 tokens/s decoding throughput at 25W power consumption. In contrast, TeLLMe delivers 25 tokens/s at only 4.8W, resulting 1.9$\times$ better decoding energy efficiency despite the GPU's 4$\times$ higher memory bandwidth, highlighting the benefits of BitNet's reduced memory access requirements.

The Raspberry Pi 5, representing a CPU-based edge solution, operates at 7.8W and achieves 16.6 tokens/s decoding throughput. TeLLMe surpasses this performance with 1.5$\times$ higher throughput while consuming 38\% less power, resulting in 2.45$\times$ better decoding energy efficiency. This comparison is particularly relevant for cost-sensitive edge deployments where both bill-of-materials cost and operational power budgets are critical constraints.

When considering the Intelligence/J metric, TeLLMe achieves the highest decoding intelligence efficiency (0.407) among all evaluated platforms, representing 1.9$\times$ improvement over the Jetson Orin Nano (0.218) and 4.6$\times$ improvement over the Raspberry Pi 5 (0.089). This metric accounts for both model quality (perplexity) and energy consumption, providing a holistic measure of inference efficiency. The superior Intelligence/J ratio validates that TeLLMe's combination of BitNet's compact representation and FPGA-optimized architecture delivers excellent inference quality per unit of energy, making it well-suited for edge applications where battery life and thermal constraints are paramount.

Furthermore, TeLLMe's perplexity (12.79) is competitive with larger models like TinyLLaMA-1.1B-W4A16 (12.42) while operating at significantly lower power. This demonstrates that BitNet's aggressive quantization maintains model quality effectively, enabling efficient edge deployment without substantial accuracy degradation. The comparison with LLaMAF's perplexity (8.89) shows room for model quality improvement, which we identify as future work through additional training sample and BitNet model architectural refinements.

\subsection{Ablation Study}
%\vspace{-1mm}
\subsubsection{Effect of TLMM Design Choice}\label{SEC:ABLATION_TLMM}
% \todo[inline]{The below part can be merged with the experiment comparison part, and further exp is required}
% \todo[inline]{We can compare our lut utilization with prior work of LUT net}
As for comparison, the LUT consumption of different TLMM unit design methods presented in \ref{sec:TLMM_METHODS} is given. The configuration for the TLMM for ablation study is set as \( G = 3 \), \( T = 28 \), and \( Q = 16 \) to guarantee the same level of parallelism and latency. Our design of the TLMM core consumes 23,082 LUTs, while the naive selection (Method 1), which selects whether to add or subtract, requires 43,176 LUTs. Another approach (Method 2) involves storing half of the possible combinations (13 out of 27) instead of all combinations and using the index to determine whether the value should be negative. This approach results in a smaller distributed RAM size, aiming to save LUTs. However, after synthesizing, it consumes 35,200 LUTs in parallel logic in $T\times G$ looking up that inverts the bit when the value is identified to be negative. The detail decompostion of the LUT usage is displayed in Table \ref{tab:LUT}, which supports the fact that $\text{LUT}_{\text{LPL}}$ is dominant in LUT consumption.






% Please add the following required packages to your document preamble:
% \usepackage{graphicx}
% \begin{table}[]
% \centering
% \caption{ablation study on weight loading unit}

% \label{tab:weight_load_ablation}
% \resizebox{\columnwidth}{!}{%
% \begin{tabular}{cccc}
% \hline
%                           & Decoding Throughput & AXI Bitwidth & Hardware Logic \\
%                           & (Token/s)           & Utilization  & Resource       \\ \hline
% 1 AXI Ports; 1:1 Transfer* & 10.3                & 54.7\%       & 1096 LUTs      \\
% 3 AXI Ports; 1:1 Transfer & +8.6                & 54.7\%       & +4289 LUTs     \\
% 3 AXI Ports; 5:3 Transfer & +12.9               & 91.1\%       & +1505 LUTs     \\ \hline
% \end{tabular}%
% }
% \footnotesize{*Baseline Configuration. m:n transfer means sending m weight data pack in n AXI transfers. AXI Interface Bitwidth = 256 bits }

% \end{table}

% \subsubsection{Effect of WBMU Design Choice} \label{sec:ABLATION_WBMU}
% As discussed in s
% \todo[inline]{Better compare the 256bit bit width utilization}
% For LLM inference, the extremely large model size necessitates fetching weights from off-chip DDR during runtime. As a result, the performance of the weight transfer unit directly impacts the end-to-end throughput, particularly in the memory-bound decoding (autoregressive) stage. This challenge is further exacerbated on embedded devices, where memory bandwidth is typically limited, making efficient weight transfer a critical design objective.

% As discussed in Section [ ], each weight data frame has a size of 140 bits, and transfers between the FPGA and DDR are performed via the AXI interface. In the baseline design, a single HP port is used, with one weight frame transmitted per transfer. Table [ ] summarizes the profiling results under different configurations. The baseline achieves only ~4 GB/s of effective bandwidth, far below the theoretical peak of 19.2 GB/s on the KV260 platform. To improve throughput, we increase the number of allocated HP ports to enable parallel weight transfers. The results show that using three HP ports yields the highest performance of 11.7 GB/s, while employing all four available HP ports actually degrades performance due to data conflicts and congestion, as HP1 and HP2 share a single QoS channel.

% Moreover, since each AXI transaction supports a width of 256 bits, transmitting a single 140-bit weight frame per transaction wastes a significant portion of available bandwidth. To address this, we concatenate five weight frames into a larger composite frame and partition it into three transfers. This strategy achieves up to 91% bandwidth utilization, calculated as (140 × 5) / (256 × 3).
% \subsubsection{Effect of Fusion Operator Design Choice}
\subsubsection{Effect of Prefill Attention Design Choice}\label{SEC:ABLATION_ATTENTION}
To demonstrate the efficacy of the reverse attention unit, we implemented a baseline naive attention design as depicted in Fig.~\ref{fig:naive_schedule}. For a fair comparison, both designs employed an identical number of processing elements, with $N_{\text{pe}} = 8$. Our reverse attention design utilized 115 DSPs, compared to 105 for the naive implementation, owing to the additional exponential operations. However, in on-board latency measurements at a prefill sequence length of 128, our design achieved a latency of 7.6 ms, whereas the naive attention required 14.3 ms. This performance improvement validates the advantages of our module on bandwidth-constrained edge FPGAs. 
%\todo[inline]{ADD the relevant BRAM}
% \subsubsection{Effect of LM Head Design Choice} 
% \todo[inline]{ADD the relevant BRAM}




\section{Conclusion}
%\vspace{-1mm}
In this work, we presented TeLLMe, the first end-to-end FPGA accelerator for ternary LLMs. TeLLMe accelerates both prefill and decoding stages. 
We propose to integrate the table-lookup-based matrix multiplication into the TeLLMe accelerator, achieving up to 25 tokens/s generation throughput and up to 143 tokens/s prefill throughput while consuming under 5W power budget and maintaining great model performance. TeLLMe achieves a significant performance and efficiency improvement over existing mobile-edge devices and previous FPGA-based accelerators, 
%% establishing a new baseline for 
enabling energy-efficient, low-latency generative AI in the edge environment.

% We proposed a table-lookup-based ternary matrix multiplication, an efficient and resource-saving matrix multiplication unit, which reused grouped activations and online precomputations for ternary matrix multiplications across projection and feedforward layers. Additionally, we introduced a fused attention unit for prefill and decoding, incorporating a reversed attention mechanism and Flash Attention-like kernel fusion. This design reduced AXI bus bandwidth usage, avoided redundant masked computation, and ensured parallelism. Furthermore, we achieved up to 9.51 tokens per second in generation, with support for 1024-token contexts, all while maintaining power consumption below 7 watts, outperforming mobile SoCs with much lower power budgets.
% TeLLMe achieved a prefill latency from 0.55s to 1.15s for prompt sizes of 64 to 128 tokens and delivered up to 9 tokens per second in decoding throughput, supporting 1024-token context lengths on edge FPGAs, with power consumption under 7 watts. This represented a significant advancement over existing mobile-edge devices and previous FPGA-based accelerators, which typically required higher precision and greater power budgets.

% To the best of our knowledge, TeLLMe was the first accelerator to provide end-to-end support for ternary LLM inference—covering both prefill and decoding stages—on real FPGA hardware, establishing a new baseline for energy-efficient, low-latency generative AI at the edge.
%%
%% The next two lines define the bibliography style to be used, and
%% the bibliography file.
\bibliographystyle{unsrt}
\bibliography{references}


%%
%% If your work has an appendix, this is the place to put it.
\appendix


\end{document}
\endinput
%%
%% End of file `sample-sigconf-authordraft.tex'.
